\documentclass[11pt]{article}

\usepackage[margin=1in]{geometry}
\usepackage{amsmath,amssymb}
\usepackage{booktabs}
\usepackage{listings}
\usepackage{xcolor}
\usepackage[hidelinks]{hyperref}
\usepackage{float}                     % [H] placement: pin code-listing figures where written
\usepackage{tabularx}                  % auto-wrapping table columns (compliance-packs table)
\usepackage{microtype}                 % better justification; reduces overflow into the margin
\setlength{\emergencystretch}{4em}     % last-resort stretch so long code tokens don't run off the page

% --- IML syntax highlighting (mirrors desktop/src/iml/iml-lang.ts: IML = a pure subset of
%     OCaml, plus the five verification directives and the ==> / <==> logical operators) ---
\definecolor{imlkw}{rgb}{0.13,0.29,0.53}      % OCaml keywords (let/rec/match/type/...)
\definecolor{imldir}{rgb}{0.55,0.11,0.47}     % IML directives (theorem/lemma/verify/instance/axiom)
\definecolor{imltype}{rgb}{0.10,0.44,0.37}    % built-in types + capitalized constructors
\definecolor{imlcomment}{rgb}{0.45,0.45,0.45} % (* ... *)
\definecolor{imlstr}{rgb}{0.00,0.40,0.15}     % "strings"
\definecolor{imlop}{rgb}{0.66,0.14,0.14}      % logical operators ==> <==>
\definecolor{imlattr}{rgb}{0.60,0.40,0.10}    % Imandra attributes [@@by ...] [@@@import ...]

\lstset{
  basicstyle=\small\ttfamily,
  columns=fullflexible,
  keepspaces=true,
  breaklines=true,
  frame=single,
  framesep=5pt,
  xleftmargin=6pt,
  aboveskip=8pt, belowskip=8pt,
  keywordstyle=[1]\color{imlkw}\bfseries,
  keywordstyle=[2]\color{imldir}\bfseries,
  keywordstyle=[3]\color{imltype},
  commentstyle=\color{imlcomment}\itshape,
  stringstyle=\color{imlstr},
}

\lstdefinelanguage{IML}{
  alsoletter={_},
  sensitive=true,
  morekeywords=[1]{and,as,assert,asr,begin,class,constraint,do,done,downto,else,end,%
    exception,external,for,fun,function,functor,if,in,include,inherit,initializer,land,%
    lazy,let,lor,lsl,lsr,lxor,match,method,mod,module,mutable,new,object,of,open,or,%
    private,rec,sig,struct,then,to,try,type,val,virtual,when,while,with},
  morekeywords=[2]{verify,instance,theorem,lemma,axiom},
  morekeywords=[3]{bool,int,int32,int64,float,char,string,unit,list,array,option,result,%
    ref,exn,bytes,nativeint,Some,None,VResult,ReuseId,MintNew,CarriedForward,%
    Fresh,Stale,Gone,Unknown,Detached},
  morecomment=[s]{(*}{*)},
  morestring=[b]",
  % Imandra attributes: [@ ...], [@@by ...], [@@@import ...] (all begin with "[@")
  moredelim=[s][\color{imlattr}]{[@}{]},
  literate={<==>}{{{\color{imlop}<==>}}}4 {==>}{{{\color{imlop}==>}}}3,
}

% The offline ponens CLI transcripts: highlight the command word and shell-style comments.
\lstdefinelanguage{PonensCLI}{
  alsoletter={_},
  sensitive=true,
  morekeywords=[1]{ponens},
  morecomment=[l]{\#},
}

% --- Pseudocode: control keywords in bold, classification outcomes in the type colour,
%     "<-" rendered as the assignment arrow; comments use the shared grey style. ---
\lstdefinelanguage{Pseudo}{
  alsoletter={_},
  sensitive=true,
  morekeywords=[1]{for,each,in,if,else,return,every,is,empty,and,or,not,intersect},
  morekeywords=[2]{CarriedForward,NeedsRereasoning},
  keywordstyle=[1]\bfseries,
  keywordstyle=[2]\color{imltype}\bfseries,
  morecomment=[l]{\#},
  literate={<-}{{$\leftarrow$}}2 {|}{{$\mid$}}1,
}

% --- The ponens policy grammar (mirrors website/src/lib/formula-fmt.mjs classifyWord/SYM_TIP):
%     temporal/lineage operators, verification statuses, CamelCase action types, and the logical
%     connectives, which the ASCII surface (-> /\ \/) renders as their canonical symbols. ---
\definecolor{polop}{rgb}{0.42,0.27,0.75}   % temporal / lineage operators (t-op)
\definecolor{poltype}{rgb}{0.20,0.40,0.75} % action types, CamelCase atoms (t-type)
\definecolor{polkw}{rgb}{0.18,0.55,0.28}   % verification statuses (t-kw)
\definecolor{polsym}{rgb}{0.42,0.45,0.50}  % logical / relational symbols (t-sym)
\lstdefinelanguage{Policy}{
  alsoletter={_},
  sensitive=true,
  morekeywords=[1]{G,F,X,P,H,U,S,S_last,P_target,P_chain,H_target,H_chain,count},
  morekeywords=[2]{completed,failed,proved,refuted,sat,unknown},
  morekeywords=[3]{EditFile,CreateFile,DeleteFile,ReadFile,SearchCode,AnalyzeCode,%
    ReadDocumentation,GitCommit,GitPush,Verify,Formalize,Decompose,GenerateTests,DefineVG,%
    Deploy,StaticAnalysis,CoverageAnalysis,Lint,TypeCheck,Decision,Output,Recommendation,%
    Plan,Draft,UserApproval,Change,Incident,Release,ToolCall,Retrieve,Compute,RunTests,%
    RunCommand,VerificationResult,Decomposition,RegionDecomposition,IMLModel,%
    CertificationCase,HLR,LLR,SourceCode,SystemRequirement},
  keywordstyle=[1]\color{polop}\bfseries,
  keywordstyle=[2]\color{polkw},
  keywordstyle=[3]\color{poltype},
  morecomment=[l]{\#},
  % ASCII surface -> canonical symbols (t-sym gray): implication, conjunction, disjunction, negation
  literate={->}{{\color{polsym}$\rightarrow$}}1 {/\\}{{\color{polsym}$\wedge$}}1
    {\\/}{{\color{polsym}$\vee$}}1 {~}{{\color{polsym}$\neg$}}1,
}

\newcommand{\trace}{\ensuremath{\sigma}}
% \code renders inline monospace and permits line breaks at natural identifier
% seams (underscores), so long snake_case tokens wrap instead of running into the margin.
\newcommand{\codeus}{\textunderscore\penalty100\relax}
\newcommand{\code}[1]{\texttt{\let\_\codeus #1}}
% --- math notation for types/records/formulas/rules (the formal-object figures) ---
\newcommand{\id}[1]{\ensuremath{\mathsf{#1}}}   % identifier: type / field / constructor / predicate / literal (works in text or math)
\newcommand{\opt}[1]{{#1}^{?}}            % optional field type
\newcommand{\many}[1]{{#1}^{*}}           % list / repeated field type

\title{Ponens: An Open-Source Framework for Verifiable, Auditable, and Composable Agent Reasoning\\[3pt]
\large The Reasoning Trace as a Defeasible, Proof-Carrying, Composable Context}
\author{Denis Ignatovich\,\textsuperscript{1} \qquad Grant Passmore\,\textsuperscript{1,2}\\[3pt]
\normalsize \texttt{denis@imandra.ai} \qquad \texttt{grant@imandra.ai}}
\date{}

\begin{document}
\maketitle
\footnotetext[1]{Imandra, Inc.}
\footnotetext[2]{University of Cambridge}
\setcounter{footnote}{2}

\begin{abstract}
Agentic AI produces reasoning at scale, but that reasoning is ephemeral, non-reproducible, and --- most
consequentially --- \emph{not composable}: there is no account of what a reasoning result depends on, so no
principled way to say what a later change invalidates. Audit logs record \emph{what happened}; they are not
reasoning \emph{states}. We present \textbf{Ponens}, an open-source framework that makes an agent's reasoning
trace a first-class logical object: a \textbf{defeasible, proof-carrying context} $\Gamma$ in which every
claim carries a certificate --- which oracle produced it, under what assumptions, whether it is \emph{fresh},
and whether it is \emph{contested}.
The trace is an append-only \emph{state} with well-formedness invariants, simultaneously \emph{memory}, a
\emph{governance surface}, and a \emph{logical context}. Over it an
\emph{evidence logic} judges an agent's work on two axes --- \textbf{goals} (the definition of done, as
outcome properties of what was established) and \textbf{policies} (temporal constraints on how it was
established, run by an \textrm{LTL}-style model checker over the action timeline) --- both resolved against
the trace rather than self-reported. Because the trace is content-addressed, its verdicts are re-checkable by
any party and any edit is tamper-evident; cryptographic sign-offs over the content hash (SSH, GPG, or keyless
sigstore), each bearing the signer's \emph{role} and \emph{disposition}, turn it into a non-repudiable
\textbf{audit trail} in which the party that did the work does not grade it. An
important technical contribution is \textbf{composition over traces}: a merge is the composition operator of
the evidence logic, and we show how to decide \emph{soundly} what a merge invalidates versus what it
provably preserves --- the moment at which two individually-valid results can become jointly invalid. We
give a durable notion of component identity that lets evidence bind to code across renames and moves. The
\emph{entire framework} is given a formal model --- a single state-transition machine in which the trace is
the state --- and machine-checked in
ImandraX\footnote{ImandraX is Imandra's automated reasoning engine: \url{https://www.imandrax.dev}.}: every
invariant, both evidence-logic axes, and the merge composition are discharged over that one state, the system
verifying its own logic. The implementation is reasoner-agnostic at its core, with a concrete instantiation
over IML/ImandraX and a \code{git} post-merge trigger that reports, on a real merge, exactly which
established results the merge put at risk. Finally we argue the same construction underpins \emph{multi-reasoner} composition: a cross-branch
dependency and a cross-reasoner dependency are the same object --- an explicit, tracked assumption.
\end{abstract}

{\setcounter{tocdepth}{2}\tableofcontents}

\part{The reasoning trace}

\section{Introduction}

An agent that writes code, proves a property, generates tests, or discharges a compliance obligation
produces \emph{reasoning}. Today that reasoning is treated as exhaust: it scrolls past in a transcript, or
is flattened into an audit log, and is gone. Two problems follow. First, it is not \emph{reproducible} ---
a reader cannot re-derive the result. Second, and more fundamentally, it is not \emph{composable}: nothing
records what a result depended on, so when the code changes --- an edit, a pull, a branch merge --- there is
no principled way to determine which prior conclusions still hold. Testing and review sample the space; they
cannot certify that a change left a body of reasoning intact, and they are structurally blind to the case
where two \emph{individually valid} changes are \emph{jointly} invalid. This paper presents the foundations
of \emph{Ponens}, an open-source framework\footnote{\code{ponens} is open source: \url{https://www.ponens.dev}
(documentation and the full trace, goal, and policy specifications) and
\url{https://github.com/imandra-ai/ponens} (source).} in which an agent's reasoning is not exhaust but a
first-class, composable, verifiable record.

Our thesis is that the right object is not a log but a \textbf{state}: the reasoning trace should be a
first-class logical context that the system computes \emph{over}, in which every recorded claim carries the
warrant that makes it usable as a premise. We make this precise and show it is not merely an accounting
convenience: it is what makes reasoning \emph{compose}.

\subsection{Motivating examples}\label{sec:intro-examples}

Five failures, all but invisible to testing and review as practiced today, motivate the construction. Each is
a case where individually-reasonable work yields a silently-wrong conclusion because nothing recorded what
that conclusion \emph{depended on}.

\paragraph{1. The merge that silently breaks a proof.} On one branch an agent proves that \code{charge} never
double-charges on a retry. On another, a colleague refactors \code{apply\_discount} --- a helper in that
proof's dependency closure --- to change rounding. Both branches pass their own tests and reviews; each result
is valid in isolation. The merge yields code in which the proof no longer holds, yet CI stays green because no
test exercises the retry-and-discount interaction. The reasoning was \emph{individually valid but jointly
invalid}, and nothing raised a signal at the merge. Deciding \emph{soundly} which established results a merge
preserves and which it puts at risk is composition over traces (\S\ref{sec:merge}).

\paragraph{2. Stale evidence that still reads as green.} A verification result from last week sits in the
record with a green checkmark. Since then a function in its dependency closure changed, but nothing
re-examined the proof --- so the dashboard still shows it passing: evidence that is stale yet presented as
current. Freshness must be a \emph{derived} verdict recomputed from the current code, never a stored flag
(\S\ref{sec:state}, \S\ref{sec:format}); a claim resting on stale evidence is not ``done.''

\paragraph{3. ``Done'' that was never established.} An agent writes ``I verified the balance invariant holds''
in its summary. There is no proof, no test --- only prose --- and a reviewer reading a plausible transcript
takes it at face value. The gap between what was \emph{asserted} and what was \emph{established} is where most
agent-work failures hide. A goal resolves only against a checkable artifact in the component's lineage, never
from prose (\S\ref{sec:goals}); an unbacked ``verified'' is recorded as \code{Unverified}.

\paragraph{4. A rename that orphans the evidence.} A refactor renames \code{charge} to \code{charge\_customer}.
Every proof, test, and decomposition about it is now attached to a name that no longer exists; naively, the
evidence is lost and the ``new'' function looks entirely unverified. Durable component identity
(\S\ref{sec:identity}) binds evidence to a code element across renames and moves, so a rename is one component,
not a delete plus an add.

\paragraph{5. Compliant only in hindsight.} A change ships correct but was produced in a way that violates a
required process --- committed before the tests ran, or merged without the independent review a regulated
codebase demands. The artifact is right; \emph{how} it was produced is not, and a flat log of what happened
cannot answer a temporal question about ordering. Policies (\S\ref{sec:policies}) are temporal constraints on
the action timeline, evaluated against the trace, so ``done'' and ``done in a compliant way'' are separate,
checkable verdicts.

None of these is exotic; together they are the ordinary texture of multi-agent, multi-branch development. What
they share is a single missing ingredient --- a record of \emph{what each conclusion depended on}, kept in a
form the system can compute over. That record is the reasoning trace.

\paragraph{Contributions.}
\begin{itemize}
\item A formal model of the reasoning trace as state (\S\ref{sec:state}, \S\ref{sec:format}) and the
argument that it is a \emph{defeasible, proof-carrying context}, with three properties --- non-monotonic,
query-carrying, and well-founded --- that make it a \emph{good} context (\S\ref{sec:context}).
\item \textbf{An evidence logic in two axes}: \emph{goals}
(\S\ref{sec:goals}) --- a property language of typed, lineage-resolved acceptance criteria over code
components --- and \emph{policies} (\S\ref{sec:policies}) --- a first-order temporal language ($G$/$F$/$U$/$S$)
over the action timeline; both resolved against the trace, with worked examples.
\item \textbf{Composition over traces} (\S\ref{sec:merge}): a merge as the composition operator of the
evidence logic, with a sound \emph{affected-set} decision (carry forward the provably-unaffected, flag the
rest) and a checkable \emph{totality} invariant.
\item \textbf{Durable component identity} (\S\ref{sec:identity}): binding evidence to code across
rename/move so composition survives refactoring, with a never-conflate resolver.
\item \textbf{A formal model and its verification} (\S\ref{sec:proved}): the framework modelled as a single
state-transition machine (the trace as state), its entire invariant set proved by the system on itself with
ImandraX --- 195 proof obligations over one state, 0 failures.
\item An implementation (\S\ref{sec:impl}), its interchange / audit / review / compliance surfaces
(\S\ref{sec:interchange}), and the \emph{multi-reasoner} generalization (\S\ref{sec:multireasoner}).
\end{itemize}

\subsection{Overview: an evidence logic over the trace}\label{sec:evidence-logic}

An agent's work is judged over the trace by an \emph{evidence logic} with two axes. \textbf{Goals} state
\emph{what} had to be established --- outcome properties, resolved against the evidence recorded in the trace
(\S\ref{sec:goals}); \textbf{policies} state \emph{how} the work had to be produced --- temporal constraints
over the timeline of actions (\S\ref{sec:policies}). These are one \emph{two-sorted} logic: the same
connectives and quantifiers range over an \emph{artifact} sort (for goals) or an \emph{action} sort (for
policies). Their verdicts --- \code{met} and \code{governed} --- are \emph{orthogonal by construction} (how
the work was done cannot change what was established), a property we machine-check (\S\ref{sec:proved}). A
third axis, \code{certified} --- a sign-off by a party \emph{other than} the producer --- rides the same
structure, which is what lets the trace be verified across parties. Because every claim carries its warrant,
all three axes are decided by \emph{re-checking the trace}, never by trusting a report; and at a merge,
composition (\S\ref{sec:merge}) re-evaluates them all.

\section{The reasoning trace as state and context}\label{sec:state}

\subsection{State, operations, and invariants}

The state \trace{} is an \textbf{append-only} record with five parts: an ordered log of \emph{actions}; a
lineage DAG of typed \emph{artifacts} (the evidence --- models, proofs, decompositions, tests, edits), each
produced by an action and linked to its inputs by \code{derived\_from}; a \emph{goal} lens (the definition
of done); a \emph{residual surface} (the negative space --- assumptions, gaps, counter-evidence); and a
\emph{derived} freshness verdict over each result. Two kinds of operation act on \trace{}:
\[
\text{\code{extend}}(\trace, op) \to \trace' \sqsupseteq \trace
\qquad
\text{\code{query}}(\trace, op) \to \text{answer} \quad (\trace \text{ unchanged}).
\]
\code{extend} appends established evidence, declares a residual, or sets the goal; \code{query} resolves a
goal, checks a policy, or renders a view. Freshness is never stored --- it is recomputed from the current
code, so the state cannot silently disagree with reality.

\paragraph{Two logics.} The trace separates the \emph{object logic} --- where a domain's behaviour is
analysed and results are produced (in our instantiation, IML discharged by the ImandraX
engine~\cite{imandra}) --- from the \emph{evidence logic}: properties \emph{of the trace itself}
(``there exists a fresh, uncontested proof, rooting in this component, that the property holds''). A goal is
a formula in the evidence logic; its atoms (\code{proved($c$)}, \code{fresh($e$)}, \code{conforms($S$)},
\ldots) quantify over artifacts, lineage, and freshness. This separation is what lets one substrate host
many object-level reasoners (\S\ref{sec:multireasoner}).

\paragraph{Invariants.} The state maintains five invariants, on which everything downstream rests:
\begin{description}
\item[I1 Append-only.] Evidence is immutable; re-analysis \emph{supersedes} (a new revision), never mutates.
\item[I2 Lineage-closed.] Every \code{derived\_from} edge points at a strictly \emph{earlier} result ---
the acyclicity condition that forbids circular premises.
\item[I3 Evidence-grounded.] Resolution reads only from evidence in lineage, never from prose.
\item[I4 Freshness-sound.] A result is \code{fresh} iff its target's dependency-closure content is
unchanged.
\item[I5 Conditional growth.] Re-establishing a fresh result is a no-op; only new work appends.
\end{description}
These are not asserted informally: \S\ref{sec:proved} discharges them as machine-checked theorems.

\subsection{A context is more than a log}\label{sec:context}

An audit log is written \emph{after} the fact and read by humans. A reasoning context is computed
\emph{over}. The difference is that every recorded claim in \trace{} carries a \textbf{certificate}: which
oracle produced it (a sound engine versus a heuristic model), under what \emph{assumptions}, whether it is
\emph{fresh}, and whether it is \emph{contested} by an open defeater. Reading \trace{} is therefore reading
a body of evidence with its warrant attached --- exactly the metadata a downstream step needs to decide
whether a prior result is an admissible premise. The trace serves as a context in two senses: an
\emph{evidence-logic} context (the model that goal resolution and policy checking evaluate against), and an
\emph{object-logic} context (an artifact can be \emph{projected} into a downstream reasoner's logic as a
premise, carrying its validity conditions as guards).

Three properties make it a \emph{good} context, and each is one of the invariants of \S\ref{sec:state}:
\begin{itemize}
\item \textbf{Non-monotonic} (I4). Freshness decays and defeaters retract, so \emph{adding} evidence can
\emph{withdraw} a prior conclusion. A proof that was sound goes \code{at\_risk} when its source drifts. This
is the correct model of a living codebase --- and it matches the semantics of the non-monotonic reasoners
(defeasible logic, answer-set programming, the event calculus) one might plug in at the object level.
\item \textbf{Carries the query side.} Residuals are \emph{open goals}: the negative space (what is not yet
established) is first-class alongside the positive. A context that carries its own open questions is directly
usable by goal-directed and abductive reasoners.
\item \textbf{Well-founded by construction} (I2). Lineage acyclicity means no claim can rest, transitively,
on itself: the context can always be evaluated bottom-up.
\end{itemize}
So one append-only structure is at once \textbf{memory} (what was done), a \textbf{governance surface} (what
policy checks), and a \textbf{logical context} (what may be assumed, and under what warrant). The rest of
the paper is what this buys: composability.

\section{The trace format}\label{sec:format}

Section~\ref{sec:state} gave the trace as an abstract state \trace{}: an append-only log, a
lineage DAG of typed artifacts, a goal lens, a residual surface, and a derived freshness verdict.
This section makes that model concrete. The format is the load-bearing artifact of the whole
construction: the invariants of \S\ref{sec:state}, the certificate of \S\ref{sec:context}, and the
composition of \S\ref{sec:merge} are all statements \emph{about the typed objects defined here}. We
present the canonical types faithfully but distilled, and close with a worked interchange trace that
exhibits every discriminator the rest of the paper relies on.

\subsection{Canonical model and interchange projection}\label{sec:format-canonical}

The specification is layered. The \textbf{canonical model} is the semantic source of truth: strongly
typed, algebraic where a set of cases is closed, and the model against which invariants and semantics
are defined. The \textbf{interchange model} is a JSON projection of it, for persistence, transport,
and cross-language interoperability. The design rule is one-directional --- \emph{specify the
strongest semantic model first; derive the wire format from it, not the reverse} --- so the JSON is
never authoritative. This matters for verifiability: the invariants of \S\ref{sec:proved} are proved
over the canonical types, and the wire format inherits them only because it is a projection.

The canonical trace is a typed record whose fields are exactly the five parts of \S\ref{sec:state}
plus reproducibility and governance surfaces: an ordered \code{actions} list, a \code{meta\_actions}
overlay, an \code{artifacts} lineage DAG, \code{policies}/\code{policy\_evaluations},
\code{execution\_environments}, \code{goals}, and inter-trace \code{trace\_links}. Every algebraic
variant projects to JSON under an \textbf{explicit discriminator}: artifacts carry
\code{artifact\_type}, actions carry \code{category} and \code{type}, and verification-result variants
carry tags such as \code{proved}/\code{refuted}/\code{sat}/\code{unknown}. Decoding must reconstruct
the strict model and \emph{reject} an invalid discriminator/payload combination --- the wire boundary
is where ill-typed traces are caught, not admitted.

\subsection{Actions: the ground-truth log}\label{sec:format-actions}

Actions are the ordered, atomic steps --- one entry per tool call --- and they do not exchange
free-form names; they exchange \emph{artifact identities} via \code{inputs} and \code{outputs}. Every
action shares a common record (Figure~\ref{fig:fmt-action-common}).

\begin{figure}[H]
\centering
\[
\id{action\_common} \;::=\; \Big(\;
\begin{array}{@{}l@{\;:\;}l@{}}
\id{id}              & \id{Int},\\
\id{label}           & \id{String},\\
\id{rationale}       & \id{String},\\
\id{detail}          & \opt{\id{String}},\\
\id{inputs}          & \many{\id{ArtifactId}},\\
\id{outputs}         & \many{\id{ArtifactId}},\\
\id{evidence}        & \many{\id{Evidence}},\\
\id{observations}    & \many{\id{Observation}},\\
\id{execution}       & \opt{\id{ExecutionMetadata}},\\
\id{reproducibility} & \opt{\id{ActionReproducibility}},\\
\id{meta\_action\_id} & \opt{\id{String}}
\end{array}
\;\Big)
\]
\caption{The common record shared by every action: identity, labels, the \code{inputs}/\code{outputs} artifact-id edges, evidence and observations, and optional execution and reproducibility metadata.}
\label{fig:fmt-action-common}
\end{figure}

The strict action type partitions steps into four \textbf{categories}, each with its own closed set
of types: \emph{activity} (\code{ReadFile}, \code{EditFile}, \code{RunTests}, \code{GitCommit},
\code{FormulatePlan}, \ldots), \emph{gateway} (decisions, splits, loops --- carrying the options and
the reason each was chosen or rejected), \emph{reasoning} (\code{Formalize},
\code{DefineVerificationGoal}, \code{Verify}, \code{StateSpaceAnalysis}, \code{ConformanceCheck},
\code{CoSimulate}, \code{GenerateTests}), and \emph{governance} (\code{EvaluatePolicy},
\code{RequestApproval}, \code{RecordAudit}, \code{LinkTrace}, \ldots). The category/type split is the
semantic layer; the concrete engine is recorded separately in \code{execution\_metadata}
(\code{tool = "imandrax"}, \code{method\_ = "region\_decomposition"}, a \code{determinism} tag), which
keeps the format \emph{reasoner-agnostic} --- the same \code{StateSpaceAnalysis} action denotes the
same reasoning operation whether ImandraX or another engine realized it.

Over this atomic log sits an optional \textbf{meta-action} overlay: an ordered, non-overlapping
grouping of actions into units of \emph{intent} (a goal-to-steps hierarchy via \code{parent\_id}),
carrying a \code{title}, \code{intent}, \code{outcome}, and a \code{source} on a fidelity ladder
(\code{PlanDeclared} $>$ \code{TurnSegmented} $>$ \code{IntentInferred}). This is the \emph{curated
narrative} of a session: the atomic actions remain the evidence, and the overlay is a producer's
claim about how they group --- reviewable as a unit before drilling in.

\subsection{Artifacts and lineage}\label{sec:format-artifacts}

Artifacts are the evidence. Every artifact carries a common record (Figure~\ref{fig:fmt-artifact-common}) and then, for the structured
kinds, a typed payload aligned with its constructor.

\begin{figure}[H]
\centering
\small
\[
\id{artifact\_common} \;::=\; \Big(\;
\begin{array}{@{}l@{\;:\;}l@{}}
\id{artifact\_id}       & \id{String},\\
\id{artifact\_role}     & \opt{\id{ArtifactRole}},\\
\id{name}               & \opt{\id{String}},\\
\id{format}             & \opt{\id{String}},\\
\id{revision}           & \opt{\id{Int}},\\
\id{producer\_action\_id} & \opt{\id{Int}},\\
\id{derived\_from}      & \many{\id{ArtifactId}},\\
\id{supersedes}         & \opt{\id{String}},\\
\id{content\_ref}       & \opt{\id{String}},\\
\id{summary}            & \opt{\id{String}},\\
\id{component\_id}      & \opt{\id{String}},\\
\id{metadata}           & \opt{\id{ArtifactMetadata}}
\end{array}
\;\Big)
\]
\caption{The common record shared by every artifact: the record identity \code{artifact\_id}, semantic role, the \code{derived\_from} lineage edges, the \code{supersedes} revision link, and the rename-stable \code{component\_id}.}
\label{fig:fmt-artifact-common}
\end{figure}

\begin{sloppypar}
The artifact type is a \textbf{strict union} so that payload presence is aligned with kind and
impossible combinations are unrepresentable: source code, plans, diffs, and commits carry only the
common record; the reasoning kinds --- \code{FormalizationArtifact}, \code{FormalModelArtifact},
\code{VerificationGoalArtifact}, \code{VerificationResultArtifact},
\code{StateSpaceAnalysisResultArtifact}, \code{ConformanceResultArtifact},
\code{CoSimulationResultArtifact}, \code{GeneratedTestsArtifact} --- pair the common record with a
typed payload. Two members are newer additions the rest of the paper leans on: the residual surface
is folded in as a \code{Residual} artifact (\S\ref{sec:format-residual}), and \textbf{1.10 adds}
\code{CarriedForwardArtifact} --- the positive record of a merge skip (\S\ref{sec:format-carried}).
\end{sloppypar}

\code{derived\_from} \emph{is} the lineage DAG: each edge names an immediate input, so an artifact's
\textbf{dependency closure} is the transitive back-walk. This is what invariant I2 of
\S\ref{sec:state} constrains --- an edge points only at a strictly earlier artifact, forbidding
circular premises --- and what goal resolution and freshness walk. \textbf{Revisioning} is
append-only (I1): an artifact is immutable once produced, so re-analysis \emph{supersedes} (a new
\code{revision}, linked back by \code{supersedes}) rather than mutating in place. ``The latest
result'' is thus a derived query over the supersession chain, never an overwrite --- which is exactly
what lets an auditor see the whole history and what lets freshness be recomputed against any point.

\subsection{Formal results}\label{sec:format-results}

The reasoning payloads carry the object-logic verdicts. A \code{VerificationGoal} states what is to
be established (\code{kind} in verify/instance/theorem/lemma/axiom, a \code{target\_symbol}, a
list of \code{property\_item}s); a \code{VerificationResult} carries the outcome as a tagged variant
--- \code{ProvedResult} (with the pretty-printed proof), \code{RefutedResult} (with a
counterexample), \code{SatResult}, or \code{UnknownResult}. The state-space, conformance, and
co-simulation payloads are structurally analogous verdicts over their own object logics. Each result may
additionally carry a \code{reasoning\_fingerprint}, the anchor for freshness treated in its own right in
\S\ref{sec:freshness}.

\subsection{The residual surface: the negative space}\label{sec:format-residual}

A trace records what was established --- the positive space --- but review and agent-to-agent handoff
need the \emph{negative} space: what was \emph{not} established. The residual surface is the uniform,
typed, queryable record of it, and \textbf{a residual is an artifact}
(\code{artifact\_type = "Residual"}): it carries its fields in \code{payload} and \emph{anchors into
the lineage DAG} through \code{derived\_from}, so a gap hangs off exactly the artifact it qualifies.
Positive and negative space share one addressable model.

\begin{figure}[H]
\centering
\small
\[
\begin{array}{@{}l@{}}
\id{residual\_kind} \;::=\; \id{Assumption} \mid \id{Unverified} \mid \id{OutOfScope} \mid \id{Limitation} \mid \id{OpenQuestion} \mid{}\\
\hphantom{\id{residual\_kind} \;::=\;{}}\id{Defeater} \mid \id{NeedsRereasoning} \mid \id{CoverageRegression}
\end{array}
\]
\begin{center}
\begin{tabular}{@{}ll@{}}
$\id{Assumption}$         & \text{\footnotesize\itshape a premise relied on but not established}\\
$\id{Unverified}$         & \text{\footnotesize\itshape something done but not checked or proved}\\
$\id{OutOfScope}$         & \text{\footnotesize\itshape deliberately not addressed}\\
$\id{Limitation}$         & \text{\footnotesize\itshape a constraint under which the results hold}\\
$\id{OpenQuestion}$       & \text{\footnotesize\itshape a decision deferred to review}\\
$\id{Defeater}$           & \text{\footnotesize\itshape counter-evidence against a claim (not a gap); carries a $\id{defeater\_kind}$}\\
$\id{NeedsRereasoning}$   & \text{\footnotesize\itshape a result a merge disturbed; re-establish against the merged code}\\
$\id{CoverageRegression}$ & \text{\footnotesize\itshape a goal whose scope gained an unproven member after a merge}
\end{tabular}
\end{center}
\[
\id{defeater\_kind} \;::=\; \id{Rebuts} \mid \id{Undermines} \mid \id{Undercuts}
\]
\begin{center}
\begin{tabular}{@{}ll@{}}
$\id{Rebuts}$     & \text{\footnotesize\itshape the claim may be false --- e.g.\ a counterexample to a proved property}\\
$\id{Undermines}$ & \text{\footnotesize\itshape the evidence is invalid --- e.g.\ the model does not match the code}\\
$\id{Undercuts}$  & \text{\footnotesize\itshape the inference is deficient --- e.g.\ ``passing tests do not establish this''}
\end{tabular}
\end{center}
\caption{The residual taxonomy (Pollock / SEI eliminative-argumentation): \code{residual\_kind} enumerates the gap and counter-evidence kinds (including the merge-specific \code{NeedsRereasoning} and \code{CoverageRegression}), and \code{defeater\_kind} classifies a defeater as a rebut, undermine, or undercut.}
\label{fig:fmt-residual-kind}
\end{figure}

The first five kinds record \emph{missing positive space} --- something not established. A
\code{Defeater} is different in kind: it records \emph{negative evidence}, a concrete reason to
believe a stated result is \emph{wrong}, and it anchors (via \code{derived\_from}/\code{target}) to
the claim it contests while citing the counter-evidence in \code{related\_artifact\_ids}. An
\emph{open} defeater makes its target \textbf{contested} --- a \code{Property} acceptance item over a
contested proof resolves \code{AcceptBlocked}, not \code{AcceptDone} (this is the non-monotonic
retraction of \S\ref{sec:context}, made first-class and reviewer-raisable). \textbf{1.10 adds} two
merge-specific kinds: \code{NeedsRereasoning} (an established result whose evidence a merge disturbed)
and \code{CoverageRegression} (a goal whose scope gained an unproven member) --- the negative half of
the composition account of \S\ref{sec:merge}. Each residual also carries a \code{severity}
(impact-if-wrong, not probability), a \code{source} on a trust ladder (\code{AgentDeclared},
\code{PolicyDerived}, \code{ToolInferred}, \code{ReviewerAdded}), a \code{status}
(\code{Open}/\code{Acknowledged}/\code{Addressed}/\code{Waived}), and an optional
\code{suggested\_check} that turns a warning into a work-item. Because it is an ordinary artifact,
every severity/status/source field and every policy applies to a defeater unchanged.

\subsection{Carried-forward: the positive dual}\label{sec:format-carried}

Composition (\S\ref{sec:merge}) needs a symmetric record on the positive side. \textbf{1.10 adds} the
\code{CarriedForwardArtifact}: where a residual records a gap, a \code{CarriedForward} records a claim
that was \emph{deliberately not re-checked because it is provably safe}. It carries a falsifiable
witness so an auditor can re-verify the skip (Figure~\ref{fig:fmt-carried-forward}).

\begin{figure}[H]
\centering
\small
\[
\id{carried\_forward\_basis} \;::=\; \id{ClosureDisjoint} \mid \id{UninterpretedOpaque}
\]
\begin{center}
\begin{tabular}{@{}ll@{}}
$\id{ClosureDisjoint}$     & \text{\footnotesize\itshape the result's dependency closure did not intersect the merge's change set}\\
$\id{UninterpretedOpaque}$ & \text{\footnotesize\itshape every touched dependency was assumed uninterpreted}
\end{tabular}
\end{center}
\[
\id{carried\_forward\_payload} \;::=\; \Big(\;
\begin{array}{@{}l@{\;:\;}l@{}}
\id{result\_id}       & \id{String},\\
\id{basis}            & \id{carried\_forward\_basis},\\
\id{closure}          & \many{\id{ComponentId}},\\
\id{via\_assumptions} & \many{\id{String}}
\end{array}
\;\Big)
\]
\caption{The \code{CarriedForward} record documenting a provably safe merge skip: the \code{basis} (disjoint closure or uninterpreted dependency) and the \code{closure} of component ids that serves as the auditor's falsifiable witness.}
\label{fig:fmt-carried-forward}
\end{figure}

The merged trace's account of every prior result is thus exactly one of three under the
\textbf{totality} invariant of \S\ref{sec:merge}: \code{CarriedForward}, \code{NeedsRereasoning}, or
freshly re-established. The two-parent \code{merge\_event} that occasions this is recorded on the
merged trace as provenance kept \emph{off} the artifact DAG --- so I2 acyclicity (a \code{derived\_from}
edge points only at an earlier artifact) survives composition. The change set is matched by
\code{component\_id}, so a rename is one component rather than a delete-plus-add; detail is deferred to
\S\ref{sec:merge}.

\subsection{Reproducibility and execution environments}\label{sec:format-repro}

Reproducibility is recorded at two granularities. At the \textbf{trace} level, a
\code{trace\_reproducibility} record states a \code{status}
(\code{NotReproducible}/\code{Partially}/\code{Reproducible}), entrypoints, and the required artifacts
and environments. At the \textbf{action} level, each consequential action may carry a
\code{reproduction\_kind} (deterministic replay, tool re-execution, procedural or manual replay), a
concrete \code{procedure} (command, working directory, arguments, or steps), and an
\code{expected\_output}. Both refer to named \code{execution\_environment}s --- a toolchain,
container, service, or external system with its versioned components --- so a reader can reconstruct
\emph{where} a result was produced, not merely that it was. The discipline is that a consequential
reasoning outcome is either directly reproducible or linked to a reproducible downstream validation.

\subsection{A worked trace}\label{sec:format-example}

Figure~\ref{fig:fmt-worked-trace} shows the interchange projection of a minimal but complete slice: one reasoning
action, the \code{VerificationGoal} it consumed and the \code{VerificationResult} it produced (with a
\code{reasoning\_fingerprint}), and one \code{Residual} contesting that result. It exhibits the
\code{category}/\code{artifact\_type} discriminators and the \code{derived\_from} lineage that
carries the whole construction.

\begin{figure}[H]
\begin{lstlisting}[language=,basicstyle=\footnotesize\ttfamily]
{
  "actions": [
    { "id": 4, "category": "reasoning", "type": "Verify",
      "label": "Verify amount_inv on the payment model",
      "rationale": "capture must never exceed the authorized amount",
      "inputs": ["a2", "g1"], "outputs": ["vr1"],
      "execution": { "tool": "imandrax", "version": "2026.04",
                     "method": "model_check", "determinism": "deterministic" } }
  ],
  "artifacts": [
    { "artifact_id": "g1", "artifact_type": "VerificationGoal",
      "producer_action_id": 3, "derived_from": ["a2"],
      "component_id": "cmp:amount_inv",
      "payload": { "goal_id": 1, "kind": "VerifyGoal",
                   "description": "amount_captured <= amount",
                   "target_artifact_id": "a2", "target_symbol": "amount_inv",
                   "properties": [ { "name": "amount_inv", "status": "Pending",
                                     "src": "amount_captured <= amount" } ] } },

    { "artifact_id": "vr1", "artifact_type": "VerificationResult",
      "artifact_role": "proof", "producer_action_id": 4,
      "derived_from": ["g1", "a2"], "component_id": "cmp:amount_inv",
      "summary": "amount_inv proved",
      "payload": { "goal_id": 1, "goal_artifact_id": "g1", "status": "proved",
                   "engine": "imandrax",
                   "result": { "proved": { "proof_pp": "...", "properties": [] } },
                   "fingerprint": {
                     "task_checksum": "sha256:7f3c...e1",   /* target + closure */
                     "task_shape": "sha256:0a9b...",
                     "target_symbol": "amount_inv",
                     "engine": "imandrax", "engine_version": "2026.04",
                     "model_artifact_id": "a2", "model_revision": 1 } } },

    { "artifact_id": "r1", "artifact_type": "Residual",
      "derived_from": ["vr1", "a13"],           /* anchors to the claim it contests */
      "summary": "amount_inv refuted under interleaved capture",
      "payload": { "kind": "defeater", "defeater_kind": "rebuts",
                   "statement": "capture can exceed amount under concurrent capture/refund",
                   "severity": "critical",
                   "target": { "target_type": "artifact", "target_id": "vr1" },
                   "related_artifact_ids": ["a13"],   /* the counterexample */
                   "source": "reviewer_added", "status": "open" } }
  ]
}
\end{lstlisting}
\caption{A worked interchange trace: a \code{Verify} reasoning action, the \code{VerificationGoal} \code{g1} and proved \code{VerificationResult} \code{vr1} (with its \code{reasoning\_fingerprint}), and an open \code{Residual} defeater \code{r1} that rebuts \code{vr1} --- exhibiting the discriminators and \code{derived\_from} lineage the rest of the paper relies on.}
\label{fig:fmt-worked-trace}
\end{figure}

Reading it as a context: \code{vr1} is a \emph{proved} result, and its fingerprint lets a consumer
recompute whether it is still \code{Fresh} against the current model. But \code{r1} is an
\emph{open} \code{Defeater} that \code{Rebuts} it (anchored by \code{derived\_from} and \code{target}),
so the claim is \emph{contested}: a \code{Property} obligation over \code{vr1} resolves
\code{AcceptBlocked}, not \code{AcceptDone}. The certificate travels with the claim --- which oracle
(\code{imandrax}), under what task (the fingerprinted closure), whether fresh, whether contested ---
which is precisely what the freshness of \S\ref{sec:proved} and the composition of \S\ref{sec:merge}
compute over.


\section{Freshness}\label{sec:freshness}

Freshness --- invariant I4 --- is the soundness signal the rest of the paper leans on. It is what lets
``the definition of done is met'' or ``this proof still holds'' be \emph{recomputed} against the current
code rather than trusted from when the work was done: a result is evidence only while the thing it was
computed over has not changed underneath it. This section makes that precise --- the fingerprint a result
carries, the derived verdict computed from it, and the formal model of that decision, proved in ImandraX.

\subsection{The fingerprint}\label{sec:freshness-fingerprint}

The anchor is one shared, result-kind-agnostic field. \textbf{Every} formal-reasoning result ---
verification, state-space, conformance, co-simulation --- may carry a \code{reasoning\_fingerprint}
(Figure~\ref{fig:fmt-fingerprint}).

\begin{figure}[H]
\centering
\small
\[
\id{reasoning\_fingerprint} \;::=\; \Big(\;
\begin{array}{@{}l@{\;:\;}l@{}}
\id{task\_checksum}    & \id{String} \quad\text{\footnotesize\itshape hash of the target $+$ its dependency closure}\\
\id{task\_shape}       & \opt{\id{String}} \quad\text{\footnotesize\itshape weaker structural hash}\\
\id{target\_symbol}    & \opt{\id{String}}\\
\id{engine}            & \opt{\id{String}}\\
\id{engine\_version}   & \opt{\id{String}}\\
\id{model\_artifact\_id} & \opt{\id{String}}\\
\id{model\_revision}   & \opt{\id{Int}}
\end{array}
\;\Big)
\]
\caption{The \code{reasoning\_fingerprint} anchors freshness (I4) in two parts: a strong \code{task\_checksum} over the target \emph{and its dependency closure} (plus a weaker \code{task\_shape}), and the \code{engine}/version and \code{model} revision that produced the result. Freshness is then a \emph{derived} verdict, obtained by recomputing this against the current model (Figure~\ref{fig:fmt-freshness}).}
\label{fig:fmt-fingerprint}
\end{figure}

The critical design point is that the fingerprint hashes the \emph{task} --- the target's dependency
closure --- and not the target's own source text. Hashing source alone is unsound in \emph{both}
directions: it \emph{misses} staleness when a result is invalidated by a change to a dependency the
target's text does not reflect, and it \emph{over-reports} staleness on comment or formatting edits
that leave the task unchanged. Closure-hashing makes a result go stale exactly when the logic it rests
on changes.

\subsection{The derived verdict}\label{sec:freshness-verdict}

Freshness is a \textbf{derived verdict}, never an authored field, recomputed by comparing the
stored fingerprint against the fingerprint of the \emph{current} model (Figure~\ref{fig:fmt-freshness}).

\begin{figure}[H]
\centering
\small
\[
\id{freshness} \;::=\; \id{Fresh} \mid \id{Stale} \mid \id{Detached}
\]
\[
\id{freshness}(\varphi) \;=\;
\begin{cases}
\id{Detached} & \text{if } \id{target}(\varphi) \notin M\\[2pt]
\id{Fresh}    & \text{if } \id{target}(\varphi) \in M \;\land\; \varphi.\id{task\_checksum} = \varphi'.\id{task\_checksum}\\
              & \quad\;\land\; \varphi.\id{engine\_version} = \varphi'.\id{engine\_version}\\[2pt]
\id{Stale}    & \text{otherwise \quad(checksum mismatch, or engine advanced)}
\end{cases}
\]
\vspace{-2pt}
\begin{minipage}{\linewidth}
{\footnotesize where $\varphi$ is the stored fingerprint and $\varphi'$ the fingerprint recomputed against the current model $M$. A \id{Detached} result is orphaned work: kept for audit, never counted as evidence.\par}
\end{minipage}
\caption{The derived \code{freshness} verdict, computed by recomputing the fingerprint against the current model and comparing: \code{Fresh} when the task checksum matches (and the engine has not advanced), \code{Stale} on a changed task or advanced engine version, and \code{Detached} when the target symbol no longer exists.}
\label{fig:fmt-freshness}
\end{figure}

A result without a fingerprint falls back to the earlier action-ordering heuristic (stale iff the
target was edited at a later action). This verdict is what makes a goal's \code{Property} acceptance
item resolve \code{AcceptDone} only when its latest result is \emph{proved, \code{Fresh}, and
uncontested} (\S\ref{sec:proved}): a proved-but-\code{Stale} or \code{Detached} result silently reopens
as a derived gap rather than continuing to read green. It is also, evaluated at a two-parent point rather
than against the working tree, the exact signal the merge operator uses to compute affected sets
(\S\ref{sec:merge}).

\subsection{The formal model}\label{sec:freshness-model}

Freshness is not merely described; its decision core is a \emph{specification-first} model in IML, proved
by ImandraX --- one of the machine-checked models of \S\ref{sec:proved}. We author the model directly in
IML and conform the implementation to it, rather than translating the code. The core is a pure function
from a fingerprint comparison to a verdict, with a worst-wins aggregate over the (possibly several)
sources a result depends on (Figure~\ref{fig:fresh-iml}).

\begin{figure}[H]
\begin{lstlisting}[language=IML]
type verdict = Fresh | Stale | Gone | Unknown   (* Gone = the paper's Detached *)

(* the per-source decision: is the file present, does the content match what was
   reasoned about, or did a logic-preserving "rescue" (e.g. a format-only edit) fire *)
let file_verdict (exists : bool) (frozen : 'h) (current : 'h) (rescue : bool) : verdict =
  if not exists       then Gone
  else if current = frozen then Fresh
  else if rescue      then Fresh
  else                     Stale

(* worst-wins over all contributing sources: gone > stale > fresh > unknown *)
let worse (a : verdict) (b : verdict) : verdict = (* ... *)
let rec combine (vs : verdict list) : verdict =
  match vs with [] -> Unknown | v :: rest -> worse v (combine rest)

(* --- invariants a silent bug would break, all proved by ImandraX --- *)

(* fresh is sound: never Fresh unless the file exists and the content matches
   (or a logic-preserving rescue justified it) *)
theorem fresh_is_sound (exists) (frozen current : 'h) (rescue) =
  file_verdict exists frozen current rescue = Fresh
    ==> exists && (current = frozen || rescue)

(* no false fresh: a Fresh aggregate hides no drifted and no missing contributor *)
theorem combine_fresh_clean (vs : verdict list) =
  combine vs = Fresh ==> not (mem Gone vs) && not (mem Stale vs)
[@@by induct ()]
\end{lstlisting}
\caption{The freshness decision core, authored specification-first in IML and proved by ImandraX (the
\emph{freshness lattice} model of \S\ref{sec:proved}, Table~\ref{tab:proofs}). \code{fresh\_is\_sound}
forbids a spurious \code{Fresh}; \code{combine\_fresh\_clean} is \emph{no-false-fresh} --- the exact
property the merge affected-set decision (\S\ref{sec:merge}) inherits. A third theorem,
\emph{revert-refreshens}, proves an edit-then-exact-revert returns to \code{Fresh}, since the verdict is
a pure function of content, not of edit history.}
\label{fig:fresh-iml}
\end{figure}

The payoff of proving \emph{this} core, rather than testing it, is that freshness is load-bearing
everywhere downstream: a single false \code{Fresh} would let a stale proof read as current evidence, and
every goal resolution (\S\ref{sec:goals}) and merge classification (\S\ref{sec:merge}) would inherit the
error. \emph{No-false-fresh} is the one property those depend on, and it is discharged once, here.


\part{Evidence logic and composition}

\section{Goals}\label{sec:goals}

\S\ref{sec:evidence-logic} introduced the two axes on which the evidence logic judges an agent's
work --- \emph{goals} (properties of what was established) and \emph{policies} (constraints on how it
was established) --- and gave the goal half in outline: a goal is a formula over \trace{} whose atoms
range over artifacts, lineage, and freshness, resolved against the evidence rather than self-reported.
This section deepens that half. We give the goal's construction, the three orthogonal verdicts it
carries, the typed acceptance criterion and the fuller property language it is a fragment of, the
lineage-based resolution rule and its non-monotonic honesty overlay, and an honest account of what is
authorable today versus what remains to be built.

\subsection{Constructing a goal}\label{sec:goals-construction}

A goal is the \emph{definition of done}, declared once, up front, and thereafter a passive lens. It is
declared through a single operation --- \code{declare\_goal} --- that records five things:
the \code{intent} (the user's words, restated), the \code{scope} (the files and symbols the goal
touches), optional \code{intent\_clauses} (the intent split into individually checkable propositions),
the \code{acceptance} items (the criteria that constitute ``done''), and optional inline
\code{policies} (the governance bar; see \S\ref{sec:policies}). The declared goal is embedded directly
into \trace{} as a first-class field (\code{goals}, canonicalizing to the empty list when no intent is
declared), so that ``definition of done'' and ``progress'' are computed by the same shared resolver
the rest of the substrate uses --- not by a private, per-tool notion. The declared goal's shape is
shown in Figure~\ref{fig:goal-goal-type}.

\begin{figure}[H]
\centering
\small
\[
\id{goal} \;::=\; \Big(\;
\begin{array}{@{}l@{\;:\;}l@{}}
\id{goal\_id}        & \id{String},\\
\id{intent}          & \id{String} \quad\text{\footnotesize\itshape (the change and why, in plain language)},\\
\id{intent\_author}  & \opt{\id{Authorship}} \quad\text{\footnotesize\itshape (who stated the intent, usually Human)},\\
\id{intent\_clauses} & \many{\id{String}} \quad\text{\footnotesize\itshape (intent split into checkable clauses)},\\
\id{scope}           & \many{\id{String}} \quad\text{\footnotesize\itshape (files / symbols the goal touches)},\\
\id{acceptance}      & \many{\id{AcceptanceItem}} \quad\text{\footnotesize\itshape (the end node: what ``done'' means)},\\
\id{policies}        & \opt{\id{GoalPolicies}} \quad\text{\footnotesize\itshape (the governance bar, \S\ref{sec:policies})},\\
\id{criteria\_review} & \opt{\id{CriteriaReview}} \quad\text{\footnotesize\itshape (the ``is it right?'' sign-off)},\\
\id{status}          & \opt{\id{GoalStatus}},\\
\id{meta\_action\_id} & \opt{\id{String}} \quad\text{\footnotesize\itshape (the meta-action pursuing this goal)}
\end{array}
\;\Big)
\]
\caption{The \code{goal} type recorded by \code{declare\_goal} and embedded in \trace{} as a first-class field, capturing intent, scope, acceptance criteria, policies, and the criteria-review sign-off.}
\label{fig:goal-goal-type}
\end{figure}

Two construction properties matter downstream, both consequences of the goal being decoupled from the
reasoning that satisfies it. First, re-declaring is \emph{idempotent}: re-stating a goal carries over
any already-resolved criterion rather than wiping a partly-proved one. Second, because the goal is a
formula over \trace{}'s evidence and \emph{not} the evidence itself, it can be attached after the fact
and is therefore \emph{portable} --- a team's definition of done is a small, versioned, human-authored
artifact that, once imported, re-resolves against \emph{your} trace with no evidence transferred. What
must line up is only the referenced components: a glob scope travels cleanly, while an exact symbol
reference needs that symbol present (or the durable identity of \S\ref{sec:identity} to rebind it
across a rename).

\subsection{Three orthogonal axes}\label{sec:goals-axes}

Where \S\ref{sec:evidence-logic} named \code{met}, \code{governed}, and \code{certified} as the
verdicts, here we make them precise. Trustworthy done is the conjunction of all three, and each is
independent and separately reported.

\begin{center}
\footnotesize
\begin{tabular}{@{}llll@{}}
\toprule
Axis & Question & Mechanism & Principal \\
\midrule
\code{met}       & Does the required \emph{artifact} exist?        & criterion $\leftrightarrow$ artifact by type + lineage & the doer's evidence \\
\code{governed}  & Was it produced compliantly?                    & the goal's policies over its cone      & org / context (a pack) \\
\code{certified} & Is the definition of done itself \emph{right}?  & criteria review by a non-doer          & a party $\neq$ the doer \\
\bottomrule
\end{tabular}
\end{center}

The axes are orthogonal \emph{by construction}, and the separation is deliberate. \code{met} asks only
whether the evidence \emph{exists}; it says nothing about quality. A refuted \code{VerificationResult}
still exists, so a criterion demanding a verification artifact for a component is met by it --- whether
that result had to be \emph{proved} rather than left refuted is a \code{governed} question, decided by
a policy over lineage, not a criterion field. This is why a goal can be \emph{met but not governed}:
the required artifact is present, but it was too weak (or produced non-compliantly) for the goal's bar
--- a state that today reads silently as ``done'' and here becomes first-class and surfaced. The
\code{certified} axis is orthogonal in a different direction: it asks not whether the formula holds but
whether it is the \emph{right} formula, and its normative rule is that the party who \emph{meets} a
goal must not be the sole party who \emph{decides its definition of done} --- so certification is a
sign-off by a principal other than the doer. Because how the work was done cannot change what was
established, \code{met} $\perp$ \code{governed} is a theorem we machine-check
(Table~\ref{tab:proofs}, \emph{governed $\perp$ met}).

\subsection{The acceptance criterion and the property language}\label{sec:goals-criterion}

An acceptance criterion, in the shipping encoding, is a typed record pairing a \emph{code component}
with the \emph{evidence artifact} that serves as its warrant (Figure~\ref{fig:goal-acceptance-criterion}).
That is all it encodes --- it does not say whether the artifact is good enough.

\begin{figure}[H]
\centering
\small
\[
\id{evidence\_req} \;::=\; \Big(\;
\begin{array}{@{}l@{\;:\;}l@{}}
\id{artifact} & \id{String} \quad\text{\footnotesize\itshape (the artifact \textsc{type} required in the component's lineage:}\\
              & \text{\footnotesize\itshape \quad $\id{VerificationResult} \mid \id{Decomp} \mid \id{Tests} \mid \id{Diff} \mid \id{Documentation} \mid \ldots$)}
\end{array}
\;\Big)
\]
\[
\id{acceptance\_criterion} \;::=\; \Big(\;
\begin{array}{@{}l@{\;:\;}l@{}}
\id{id}        & \id{String},\\
\id{statement} & \id{String} \quad\text{\footnotesize\itshape (human ``what must be true'')},\\
\id{component} & \big(\id{function\_} : \id{String},\; \id{file} : \opt{\id{String}}\big),\\
\id{evidence}  & \id{evidence\_req},\\
\id{required}  & \id{Bool},\\
\id{covers}    & \many{\id{String}} \quad\text{\footnotesize\itshape (intent clauses it addresses)}
\end{array}
\;\Big)
\]
\caption{The shipping acceptance-criterion encoding: a typed record binding a code component to the evidence-artifact type required as its warrant, together with the intent clauses it covers.}
\label{fig:goal-acceptance-criterion}
\end{figure}

This fixed struct is a \emph{concrete encoding of a fragment} of a fuller \textbf{property language} ---
the evidence logic proper --- and every legacy form desugars into it, so no existing trace breaks. The
property language expresses a goal's ``done'' as a conjunction of \emph{obligations}, each a predicate
over \trace{} tagged with the verdict axis it feeds. The atoms are exactly the resolver predicates the
implementation already computes, listed in Figure~\ref{fig:goal-atoms}.

\begin{figure}[H]
\centering
\small
\[
\begin{array}{@{}l@{\qquad}l@{}}
\id{has}(c,T)      & \text{an artifact of type $T$ roots in component $c$ (latest wins)}\\
\id{proved}(c)     & \id{has}(c,\id{VerificationResult}) \land \id{status} \in \{\id{proved},\id{sat}\}\\
\id{refuted}(c)    & \id{has}(c,\id{VerificationResult}) \land \id{status} = \id{refuted}\\
\id{conforms}(S)   & \text{a $\id{ConformanceResult}$, complete, $\id{diverged}=0$, over sources $S$}\\
\id{decomposed}(c) & \text{a decomposition roots in $c$}\\
\id{tested}(c)     & \text{a region-complete / coverage-rooted suite roots in $c$}\\
\id{no\_open\_gap}(c) & \text{no open residual touches $c$}\\
\id{fresh}(e)      & \text{evidence $e$ is closure-fresh (\S\ref{sec:state}, I4)}\\
\id{uncontested}(e)& \text{$e$ is not contested by an open defeater}
\end{array}
\]
\caption{The atoms of the property language --- the resolver predicates the implementation already computes over \trace{}, from artifact existence and verdict checks to freshness and contestation.}
\label{fig:goal-atoms}
\end{figure}

\noindent\textbf{Selectors} turn coverage into a quantifier rather than one row per symbol:
\code{component(f)}, \code{symbol(s)}, \code{module(p)}, \code{glob("payments/**")}, \code{tag(l)},
\code{scope}, and \code{high\_stakes} (the configured high-stakes glob set). An obligation is then
\code{[role]\ formula}, where \code{formula} combines atoms with \code{and}, \code{or}, \code{not},
\code{=>} and the quantifiers \code{forall $x$ in $sel$} / \code{exists $x$ in $sel$}, and
\code{role} \(\in\) \code{\{[met], [governed]\}} (default \code{[met]}). The role is not
``deliverable versus standard'': a \code{[met]} obligation is an \emph{outcome} property over the
resulting artifacts (\code{proved($c$)}, \code{conforms($S$)}, \code{no\_open\_gap($c$)}), while a
\code{[governed]} obligation is a \emph{process} constraint over how those artifacts were produced
(\code{formalize before verify}, \code{eventually reproved}). The connectives and quantifiers are
shared; the \emph{sort ranged over} --- the artifact sort versus the action sort --- \emph{is} the
goal-versus-policy distinction of \S\ref{sec:evidence-logic}. Freshness is a \emph{modality}, not a
truth condition: an obligation that holds but whose witness fails \code{fresh} is
\emph{true-but-\code{at\_risk}}, not false --- the honesty layer, made first-class in the language.

The desugaring makes one subtlety explicit. A typed \code{\{component, evidence\}} criterion is
\code{has($c$, T)} --- \emph{existence}, so a refuted result still satisfies it --- whereas the legacy
\code{property} kind is \code{proved($c$)}, which checks the verdict and so resolves a refuted result to
\code{blocked}. In the language one simply picks the atom that means what is intended, as the
desugarings in Figure~\ref{fig:goal-desugaring} show.

\begin{figure}[H]
\centering
\small
\[
\begin{array}{@{}l@{\;\rightsquigarrow\;}l@{}}
\big(\id{component}{:}\{\id{function\_}{:}f\},\; \id{evidence}{:}\{\id{artifact}{:}T\}\big) & [\id{met}]\; \id{has}(f, T)\\[2pt]
\id{kind}{:}\,\id{property}\ \{\id{symbol}\}      & [\id{met}]\; \id{proved}(\id{symbol})\\[2pt]
\id{kind}{:}\,\id{obligation}\ \{\id{policy\_id}\} & [\id{governed}]\; \id{holds}(\id{policy\_id})\\[2pt]
\id{kind}{:}\,\id{gap}\ \{\id{residual\_id}\}      & [\id{met}]\; \id{no\_open\_gap}(\id{target\text{-}of}\ \id{residual\_id})\\[2pt]
\id{kind}{:}\,\id{change}\ \{\id{symbol}\}         & [\id{met}]\; \id{has}(\id{symbol},\; \id{Diff} \lor \id{IMLModel})
\end{array}
\]
\caption{How each legacy criterion kind desugars into a role-tagged atom of the property language, so no existing trace breaks and one picks the atom that means what is intended.}
\label{fig:goal-desugaring}
\end{figure}

\subsection{Resolution by lineage}\label{sec:goals-resolution}

Resolution \emph{evaluates} the formula against the trace; nothing is authored as done. A criterion
\code{has($c$, T)} (and its verdict-checking cousins) resolves \code{met} when \trace{} holds an
artifact of type \code{T} whose lineage \textbf{roots in} the component $c$ --- read from structured
provenance (\code{target\_symbol}, \code{derived\_from}, the model's symbol list), never from a
description string. This is the fix for the ``goal never ticks'' defect: binding by a substring of the
goal's prose silently misses when the proved symbol's name and the description's wording diverge, so the
rooting is computed over lineage instead. When several artifacts match, the latest by producer action is
the evidence pointer. Concretely the resolver binds as shown in Figure~\ref{fig:goal-match}.

\begin{figure}[H]
\centering
\small
\[
\begin{array}{@{}l@{}}
\id{match}(o) \;=\; \operatorname*{argmax}_{a \in \mathcal{C}}\; \id{producer\_action\_id}(a), \quad\text{where}\\[4pt]
\mathcal{C} \;=\; \{\, a \;\mid\; \id{artifact}(a) \land \id{type}(a) = T \land \id{roots\_in\_component}(a, c) \land \lnot\,\id{contested\_by\_open\_defeater}(a) \,\}
\end{array}
\]
\caption{The resolver's binding rule: among artifacts of the required type whose lineage roots in the component and that are not contested by an open defeater, the latest by producer action becomes the evidence pointer.}
\label{fig:goal-match}
\end{figure}

\noindent then writes onto the \emph{enriched copy} (never back onto the authored goal) a \code{status}
(\code{done} on a match, \code{blocked} on a matched-but-refuted-or-contested result, \code{todo} on no
match), the matched artifact's id as \code{evidence\_ref}, and the \code{at\_risk} flag with a reason
when the witness is stale or detached.

\textbf{Component precision.} \code{roots\_in\_component} is \emph{target-precise}: an artifact that
declares its own \code{target\_symbol} (a verification goal, a decomposition, a targeted diff) is about
\emph{that} component, not every symbol the shared model happens to formalize. Only when nothing in the
lineage names a target does resolution fall back to the model's symbol list. Without this, a
decomposition of one function would spuriously appear to root in every symbol of the enclosing model.
Because the match is by component-plus-type over lineage rather than by a stored pointer, one artifact
can satisfy the same obligation across \emph{several} goals --- goals are lenses over one shared evidence
stream, not containers owning separate traces --- and a rename is followed
(\code{roots\_in\_component} accepts the source name \emph{or} the IML symbol; \S\ref{sec:identity}).

\textbf{Progress and the atom-to-helper mapping.} Progress is the fraction of required items resolved:
\(\text{progress} = |\{\text{resolved}\}| / |\{\text{total}\}| \in [0,1]\) (with an in-progress item
counting as a half), and a goal is \code{met} exactly when \emph{all} required items resolve. Each atom
delegates to a shipping resolver helper --- \code{proved}/\code{has} to the typed-lineage resolver over
\code{roots\_in\_component}, \code{no\_open\_gap} to the residual scan, \code{fresh}/\code{uncontested}
to the freshness verdict and the open-defeater check. These facts are machine-checked
(Table~\ref{tab:proofs}, \emph{goal resolution}): \emph{progress} \(\in [0,1]\),
\emph{met $\Leftrightarrow$ all-done}, and \emph{at\_risk never demotes}.

\textbf{The \code{at\_risk} honesty overlay (non-monotonic).} Because the trace is a non-monotonic
context (\S\ref{sec:context}, I4), a met obligation can drift. When the evidence a met item rests on
goes stale --- its target's dependency-closure content changed --- the item \emph{stays met} but is
flagged \code{at\_risk}, with a reason and a pointer to the derived stale-evidence residual. It never
silently demotes to unmet, and it never silently passes as clean: the two are distinguished, and a
re-established, fresh result clears the flag. A resolved obligation captures the closure checksum its
witness was verified against, so an edit trips exactly the same freshness authority that governs the
rest of the substrate. Figure~\ref{fig:goal-at-risk} traces a met obligation through a source edit and
a superseding re-proof.

\begin{figure}[H]
\centering
\small
\noindent\textit{Example (met but \id{at\_risk}).} The obligation $[\id{met}]\;\id{proved}(\id{refund})\;[\id{fresh}]$ is
resolved at three checkpoints:
\[
\begin{array}{@{}l@{\;\rightarrow\;}l@{}}
\text{A: \id{refund} proved fresh} & \id{done},\; \id{evidence\_ref} = \id{VR\#7},\; \id{at\_risk} = \id{false}\\[3pt]
\text{B: source of \id{refund} edited, no re-proof} & \id{done},\; \id{evidence\_ref} = \id{VR\#7},\; \id{at\_risk} = \id{true}\\
 & \quad\text{\footnotesize\itshape (``re-verify: source drifted''; still met, never silently demotes)}\\[3pt]
\text{C: superseding fresh proof appends} & \id{done},\; \id{evidence\_ref} = \id{VR\#12},\; \id{at\_risk} = \id{false}
\end{array}
\]
\caption{The non-monotonic \code{at\_risk} overlay: an obligation stays met when its witness goes stale after a source edit but is flagged for re-verification, and a superseding fresh proof clears the flag --- honesty without silent demotion.}
\label{fig:goal-at-risk}
\end{figure}

\subsection{Faithfulness --- the certified axis}\label{sec:goals-certified}

\code{met} and \code{governed} both sit \emph{downstream} of an unguarded seam: in the real workflow
the same agent that authors the definition of done then meets it, a textbook specification-gaming
setup. An agent can be perfectly honest at resolution --- never fabricate an artifact --- and still
``succeed'' by choosing a weak or incomplete definition of done (the ``wrong spec / vacuous proof''
problem). Faithfulness cannot be decided mechanically --- whether a formal acceptance faithfully
captures an informal intent is the formalization gap, and intent lives in a human's head --- so the
certified axis does not \emph{verify} faithfulness; it makes the seam visible, reviewed by a different
principal, coverage-checked, and hard to retrofit. Two mechanisms carry it. \textbf{Coverage:} each
obligation's \code{covers} field names the intent clauses it addresses, so an \code{intent\_clauses}
entry with no covering obligation surfaces as an uncovered-clause residual --- completeness becomes a
checkable projection rather than a judgement call. \textbf{Criteria review:} a \code{criteria\_review}
record signs off the acceptance \emph{set} --- ``this is what done means'' --- and its author
\emph{must} be a party other than the doer (\code{Human} or \code{Reviewer}, never the working
\code{Agent}). The certified verdict --- \code{faithfulness\_of} the goal --- is \code{Approved} when a
non-doer has signed off the definition \emph{and} no intent clause is left uncovered. A criterion that
\emph{bites} is evidence of faithfulness: a property refuted with a counterexample is stronger evidence
that the definition captures something real than an all-green suite of trivial edits. Certification also
sees the governance overrides: a goal certified with a formal-methods pack \emph{disabled} is a visibly
weaker claim than one certified with it enforced.

\subsection{Authorable today versus open}\label{sec:goals-status}

We are precise about what ships. The \emph{leaf atoms} of the property language already exist --- they
are the resolver helpers that compute \code{proved}, \code{has}, \code{no\_open\_gap}, freshness, and
contestation --- and the shipping authorable surface is the fixed \code{component}\,+\,\code{evidence}
struct, which is exactly the \code{exists}-over-one-component fragment those helpers evaluate. The
lineage-precise resolution rule of \S\ref{sec:goals-resolution}, the \code{at\_risk} overlay, and the
legacy-desugaring path are all shipped (in the reference resolver, both the Python and the JavaScript
implementations, held in parity by a harness). What is \emph{not yet built} is the full property-language
front-end that makes the language \emph{authorable} in its generality:

\begin{itemize}
\item a \emph{formula AST plus recursive evaluator} in place of the resolver's fixed \code{kind}-switch,
so \code{and}/\code{or}/\code{not}/\code{=>} and the role tags compose;
\item a \emph{component index} so \code{forall $f$ in glob("payments/**")}, \code{module($p$)}, and
\code{tag($l$)} can \emph{enumerate} their sets --- the one genuinely new piece of machinery, though it
indexes data (per-file source symbols, the high-stakes glob set) the trace already carries;
\item \emph{stable ids for \code{properties(S)}}, so a spec's authored verification goals are
addressable for \code{forall $p$ in properties(S)}.
\end{itemize}

\noindent Until then, goals are authored in the fixed struct (which desugars faithfully), and quantified
coverage is expressed through \code{[governed]} policies rather than \code{[met]} selectors. The design
of record is the two-sorted property language of \S\ref{sec:goals-criterion}; the evaluator for its
leaves is real; the combinator and enumeration layers are the honest outstanding work.

\section{Policies}\label{sec:policies}

Section~\ref{sec:goals} developed the \emph{met} axis: a goal is an outcome
property over the \emph{artifact} sort, resolved against the evidence in \trace{}.
This section deepens the second axis. A goal asks \emph{what was established}; a
policy asks \emph{how it was established}. Because both are formulas in the one
two-sorted evidence logic (\S\ref{sec:evidence-logic}), a policy shares the
connectives and quantifiers of a goal but ranges over the \emph{action} sort ---
the finite timeline of what the agent did --- and over the trace's structural
surfaces (lineage, residuals). The verdict it yields is \code{governed}, distinct
from \code{met} and, as we show in \S\ref{sec:pol-met-vs-governed}, orthogonal to
it. This is the \emph{governance surface} of \S\ref{sec:context} made precise:
governance you \emph{run} over the trace, not a checklist or a sign-off.

\subsection{The policy object}\label{sec:pol-object}

A policy is a machine-checkable rule over a trace, carrying its meaning as a
\emph{formula} rather than prose. The canonical object (POLICY\_SPEC\_v0\_2
\S6.1) is shown in Figure~\ref{fig:pol-object}.

\begin{figure}[H]
\centering
\small
\[
\id{policy} \;::=\; \Big(\;
\begin{array}{@{}l@{\;:\;}l@{}}
\id{policy\_id}            & \id{String},\\
\id{name}                  & \id{String},\\
\id{severity}              & \id{Policy\_severity},\\
\id{scope}                 & \id{Policy\_scope},\\
\id{kind}                  & \id{Policy\_kind},\\
\id{applies\_when}         & \opt{\id{Policy\_selector}},\\
\id{formula}               & \id{Policy\_formula},\\
\id{reference\_artifact\_id} & \opt{\id{String}}
\end{array}
\;\Big)
\]
\caption{The canonical \code{policy} record (POLICY\_SPEC\_v0\_2 \S6.1), whose \code{formula} field carries the rule's authoritative machine-checkable meaning. \code{severity} is \id{Info}$\mid$\id{Warning}$\mid$\id{ErrorSeverity}; \code{scope} is \id{Trace}$\mid$\id{Action}$\mid$\id{Artifact}$\mid$\id{Module}; \code{kind} is \id{TraceInvariant}$\mid$\id{ReasoningRequirement}$\mid$\id{LineageRequirement}$\mid\cdots$}
\label{fig:pol-object}
\end{figure}

\noindent The \code{formula} field is authoritative: if the human
\code{description} and the \code{formula} disagree, the formula wins.
\code{severity} records the consequence of failure and induces the ordering
\(\code{Info} < \code{Warning} < \code{Error}\); an \code{Error}-severity
failure is a \emph{hard} governance failure, a \code{Warning} a flagged gap.
\code{scope} fixes the level at which the rule is interpreted (whole trace, an
action, an artifact, or a module), and \code{applies\_when} is an optional
selector (by action type, artifact type/role, category, or path pattern) that
narrows the trigger. A policy may additionally name a required \code{reasoner}
from the reasoner registry, demanding not just \emph{that} a claim was verified
but \emph{by what} engine --- the difference between a claim \emph{established}
and one merely asserted.

A policy \emph{attaches} to a trace in two ways. It may be materialized from a
\emph{pack} or \emph{source} (\S\ref{sec:pol-sources}), or it may be declared
\emph{inline on a goal} --- the \code{policies} passed alongside
\code{declare\_goal}, evaluated over that goal's relevance cone. Inline policies
are the natural place to demand that a goal's evidence be \emph{good}, not merely
present (e.g.\ ``a refuted result must eventually be re-proved,''
\S\ref{sec:pol-worked}), while packs carry the shared standards a whole codebase
is held to.

\subsection{The formula language}\label{sec:pol-language}

The policy language (POLICY\_LANGUAGE\_v0\_2) is \textbf{LTLf} --- linear
temporal logic over \emph{finite} traces~\cite{ltl} --- extended with
\emph{scoped-past} operators and \emph{first-order structural predicates} over
the typed artifact graph. A trace is read as a finite sequence of actions
\(\tau = a_0, a_1, \ldots, a_n\) over a lineage DAG, and a formula is evaluated at
a position \(i\).

\paragraph{Grammar.}
The canonical abstract syntax (POLICY\_SPEC \S10) is the grammar of
Figure~\ref{fig:pol-grammar}, written in the unicode surface
(\(\to\), \(\land\), \(\lor\), \(\lnot\)) the gallery uses throughout.

\begin{figure}[H]
\centering
\small
\[
\begin{array}{r@{\;\;}c@{\;\;}l}
\varphi & ::= & \top \mid \bot \mid p \mid \lnot\varphi \mid \varphi \land \varphi \mid \varphi \lor \varphi \mid \varphi \rightarrow \varphi \mid \varphi \leftrightarrow \varphi \\
        & \mid & t \bowtie t \mid o_1\,\varphi \mid \varphi \mathbin{o_2} \varphi \mid Q\,x \in e.\;\varphi \\[3pt]
o_1 & ::= & \mathsf{G} \mid \mathsf{F} \mid \mathsf{X} \mid \mathsf{P} \mid \mathsf{H} \mid \mathsf{P}_{\mathsf{target}} \mid \mathsf{P}_{\mathsf{chain}} \mid \mathsf{H}_{\mathsf{target}} \mid \mathsf{H}_{\mathsf{chain}} \\
o_2 & ::= & \mathsf{U} \mid \mathsf{S} \mid \mathsf{S}_{\mathsf{last}} \\
Q & ::= & \forall \mid \exists \mid \exists! \\
{\bowtie} & ::= & {=} \mid {\neq} \mid {<} \mid {\leq} \mid {>} \mid {\geq} \mid {\in} \\
t & ::= & x \mid v \mid t.f \mid \mathrm{inputs}(t) \mid \mathrm{outputs}(t) \mid \mathrm{target}(t) \mid \mathrm{ancestors}(t) \mid \lvert e\rvert \\
e & ::= & \{\, t_1,\ldots,t_n \,\} \mid \{\, t : x \in \mathit{src},\ \varphi \,\}
\end{array}
\]
\vspace{-2pt}
\begin{minipage}{\linewidth}
{\footnotesize $p$: an atomic predicate on an action, artifact, status, or path. $\mathsf{G},\mathsf{F},\mathsf{X},\mathsf{P},\mathsf{H}$: globally, finally, next, previous, historically; the \textsf{target}/\textsf{chain} subscripts scope a past operator to a target or its lineage chain. $\mathsf{U},\mathsf{S}$: until, since ($\mathsf{S}_{\mathsf{last}}$: since the most recent occurrence). $t.f$: field access (e.g.\ $r.\mathrm{severity}$); $\lvert e\rvert$: cardinality; $e$ is a set literal or a lineage comprehension over the artifact graph.\par}
\end{minipage}
\caption{The abstract syntax of the \textrm{LTL}$_f$-based policy formula language (POLICY\_SPEC \S10): temporal operators, first-order quantifiers, field-access terms, and set comprehensions over the typed artifact graph. This is the unicode surface used throughout the gallery.}
\label{fig:pol-grammar}
\end{figure}

\paragraph{Temporal operators.}
The language supports \emph{both} future- and past-time operators over the
finite timeline, plus scoped variants:
\begin{itemize}
\item \textbf{Future.} \code{G}\,\(\varphi\) (globally --- \(\varphi\) holds at
every position); \code{F}\,\(\varphi\) (finally/eventually --- \(\varphi\) holds
at some later position); \code{X}\,\(\varphi\) (next); and the binary
\(\varphi\)\,\code{U}\,\(\psi\) (until --- \(\varphi\) holds at every position
until \(\psi\) does).
\item \textbf{Past.} \code{P}\,\(\varphi\) (previously --- held at some earlier
position); \code{H}\,\(\varphi\) (historically --- held at every earlier
position); \(\varphi\)\,\code{S}\,\(\psi\) (since --- \(\varphi\) has held at
every position since \(\psi\) last held); and \(\varphi\)\,\code{S\_last}\,\(\psi\)
(since-last). Past operators are natural here because most trace rules look
\emph{backwards} from a triggering event, and they admit \emph{incremental}
runtime evaluation (blocking an action \emph{before} it violates a policy).
\item \textbf{Scoped past.} \code{P\_target}/\code{H\_target} restrict the
witnessing earlier position to those sharing the current action's \emph{target}
(same file/module/artifact); \code{P\_chain}/\code{H\_chain} restrict it to
positions lying in the current action's \code{derived\_from} \emph{lineage}. This
is what distinguishes ``any prior work'' from ``\emph{relevant} prior work'':
\code{research\_before\_edit} uses \code{P\_target}, while
\code{reasoning\_required\_for\_high\_stakes} uses \code{P\_chain}.
\end{itemize}

\noindent The formal interpretation is standard finite-trace LTL:
\((\tau, i) \models \code{G}\,\varphi\) iff \((\tau, j) \models \varphi\) for all
\(j \ge i\); \((\tau, i) \models \code{F}\,\varphi\) iff \((\tau, j) \models
\varphi\) for some \(j \ge i\); \((\tau, i) \models \code{P}\,\varphi\) iff for
some \(j < i\); \((\tau, i) \models \code{H}\,\varphi\) iff for all \(j < i\); and
\((\tau, i) \models \varphi\,\code{S}\,\psi\) iff there is \(j < i\) with
\((\tau, j) \models \psi\) and \((\tau, k) \models \varphi\) for all
\(j < k < i\). The scoped operators are the same, with the witness \(j\)
restricted to the target- or lineage-matching positions.

\paragraph{Atoms, fields, and structure.}
At each position the atoms range over action types (\code{EditFile},
\code{GitCommit}, \code{Verify}, \code{DefineVG}, \code{Formalize},
\code{Decompose}, \ldots), action categories (\code{gateway}, \code{reasoning},
\code{activity}), status predicates (\code{completed}, \code{failed}), field
predicates (\code{rationale}\,\(\neq \emptyset\)), and path predicates
(\code{high\_stakes\_path}). Structural predicates read the typed graph:
artifact types (\code{VerificationResult}, \code{IMLModel},
\code{Decomposition}, \ldots), result status (\code{proved}, \code{refuted},
\code{sat}, \code{unknown}), lineage (\code{ancestors(derived\_from)}), and set
operations (\code{inputs}, \code{outputs}, \code{target}). Field access
(\code{Field}) with comparisons (\code{cmp\_op}) is what lets a policy talk about
the \emph{residual surface}: severity is \emph{totally ordered}
(\code{Info} < \code{Low} < \code{Medium} < \code{High} < \code{Critical}), so
\(r.\code{severity} \ge \code{High}\) is meaningful, while \code{kind} and
\code{status} are unordered enums compared by equality, and \(\emptyset\) denotes
a missing value.

Each policy declares its \code{language\_level} --- \code{pure\_temporal} (future/
past operators over atomic propositions), \code{scoped\_temporal} (adds the
scoped-past operators), or \code{structural} (adds quantifiers, comprehensions,
and lineage predicates) --- so the evaluator and the reviewer know how much
machinery a check requires.

\paragraph{A gallery sampler.}
The published gallery (\S\ref{sec:pol-sources}) is the best intuition for the
language: over two hundred real policies, from CI hygiene to avionics and market
governance, all in the one surface. Figure~\ref{fig:pol-gallery} shows a
representative handful --- each a live gallery entry, glossed by its own
\code{description} --- spanning the operator range: past-response (\code{P}),
liveness (\code{F}), the since-last binary (\code{S\_last}), the two scoped-past
forms (\code{P\_target} on the same target, \code{P\_chain} along lineage), and
negation over the connectives.

\begin{figure}[H]
\begin{lstlisting}[language=Policy]
# tests_before_commit: every commit follows a passing test run
G(GitCommit -> P(RunTests /\ completed))

# no_secrets_committed: a commit must never introduce a secret-bearing file
G(GitCommit -> ~P(CreateFile /\ secret_pattern))

# research_before_edit: every edit follows research on the SAME target
G(EditFile -> P_target(ReadFile \/ ReadDocumentation \/ SearchCode \/ AnalyzeCode))

# reasoning_required_for_high_stakes: high-stakes edits need proof in their lineage
G(EditFile /\ high_stakes_path -> P_chain(VerificationResult(proved \/ sat) \/ Decomposition))

# tests_pass_before_commit: the most recent test run before each commit passed
G(GitCommit -> (RunTests /\ completed) S_last GitCommit)

# refuted_results_must_be_reproved: a refutation must eventually be re-proved
G(VerificationResult(refuted) -> F(VerificationResult(proved \/ sat)))

# no_release_without_authenticated_approval: release only after a scoped approval
G(Release \/ Deploy -> P(UserApproval /\ authenticated /\ approval_scope_covers))
\end{lstlisting}
\caption{A sampler of live gallery policies (each line commented with its policy id
and \code{description}), rendered in the unicode surface of
Figure~\ref{fig:pol-grammar}. Operators are purple, action/artifact types blue,
status predicates green, and the connectives ($\rightarrow \land \lor \lnot$)
grey; lowercase atoms (\code{secret\_pattern}, \code{high\_stakes\_path},
\code{authenticated}) are the flags a trace carries.}
\label{fig:pol-gallery}
\end{figure}

\subsection{Evaluation semantics}\label{sec:pol-eval}

Evaluating a policy \(\varphi\) against a trace \(\tau\) computes the satisfaction
relation \(\tau \models \varphi\) and records a \code{policy\_evaluation} with one
of three (materially) statuses:
\begin{itemize}
\item \code{Passed} --- the formula holds over the trace;
\item \code{Failed} --- it is violated, and the evaluator reports the
\emph{witnessing} \code{violating\_action\_ids}\slash\allowbreak\code{violating\_artifact\_ids};
\item \code{NotApplicable} --- the policy's trigger never occurs (e.g.\ a commit
policy on a trace with no commit). \code{NotApplicable} results are \emph{excluded}
when a grade renormalizes compliance --- a policy that does not apply neither
passes nor fails a trace. (A fourth status, \code{UnknownEvaluation}, records the
rare case where the checker cannot decide.)
\end{itemize}

\paragraph{Over-empty semantics.}
Because collections can be empty and the quantifiers are the ordinary
finite-domain ones, the \emph{vacuous} cases are load-bearing and follow the
standard reading. A universal over an empty collection is \emph{passed}
(\(\forall x \in \emptyset.\,\varphi\) is vacuously true); an existential over an
empty collection is \emph{failed} (\(\exists x \in \emptyset.\,\varphi\) has no
witness). The temporal analogue is the same: \code{G}\,\(\varphi\) over a timeline
with no triggering position is vacuously satisfied --- which is exactly the
``vacuous case'' of the response pattern in \S\ref{sec:pol-worked}. This is why a
trigger that never fires yields \code{NotApplicable} rather than a spurious pass:
the evaluator distinguishes ``held everywhere it could'' from ``never had the
chance to fail.''

\paragraph{Fold-back onto \trace{} and the \code{governed} verdict.}
The \code{policy\_evaluation} records are \code{query} results
(\S\ref{sec:state}): they read \trace{} and do not mutate the evidence DAG --- a
policy \emph{constrains} history, it does not \emph{define} it. The evaluations
are folded back as a derived \emph{governance verdict} (cached, never evidence),
distinct from the \code{met} verdict of \S\ref{sec:goals}. A goal resolves
\code{met} from its acceptance obligations over the artifact sort; it resolves
\code{governed} from the policies over its cone: the inline goal policies
evaluated over the goal's relevance cone, plus the active pack's policies. An
\code{Error}-severity failure makes the goal \emph{not governed} even when it is
\code{met}; a \code{Warning} leaves it governed-with-a-flagged-gap. Severity thus
composes with the verdict rather than the formula.

\subsection{Policy sources and packs}\label{sec:pol-sources}

Policies come from more than one place, and where a policy \emph{lives} is a
distribution-and-trust concern, not a semantic one: a policy's identity and
meaning are unchanged by its source. A \textbf{source}
(POLICY\_SOURCES\_v0\_1) is a named, typed origin --- a \code{gallery} (a
published, hash-verified static catalog over HTTP), a \code{local} directory of
policy JSON committed in-repo, or a (future) authenticated \code{hub}. Sources
are layered like \code{.gitconfig}/\code{.npmrc} (user-global then project), the
order is the precedence for resolving unqualified references, and there is
\emph{no silent shadowing}: an ambiguous unqualified id across sources is an
error listing the qualified options, not a coin flip.

A \textbf{pack} is a named set of policies distributed by a source and
materialized together into a trace. The formal-methods pack
\code{apply-formal-methods} has two halves: a \emph{coverage} group (high-stakes
code --- money math, state machines, access control, bounds --- must be backed by
a proof or a decomposition in its lineage before it ships) and a \emph{rigor}
group (once the pipeline runs it must run properly: formalize before verify,
refuted goals get re-proved, decompositions produce tests). Which surface counts
as ``high-stakes'' is data-driven from the trace's \code{high\_stakes\_paths}
rather than hard-coded. Compliance packs (e.g.\ DO-178C, MISRA, NIST AI RMF,
ESMA/MiFID) are structurally identical --- sets of the same LTLf formulas.

When a policy is materialized, the CLI stamps its \emph{provenance} onto the
trace: \code{\{ source, id, version, hash \}}. A governed trace therefore records
exactly which policies, from which sources, at which content hashes, it was
checked against --- closing the loop \emph{source \(\to\) policy \(\to\) trace
\(\to\) sign-off} that a compliance reviewer needs. This is the \emph{governance
boundary} in operation: \code{ponens} \emph{evaluates} the policies over the
trace, while the producing reasoner only \emph{emits evidence}. The party that
did the work does not grade the work; the certificate travels with the claim and
is re-checkable by another party (\S\ref{sec:context}).

\subsection{Worked examples}\label{sec:pol-worked}

\paragraph{A response (temporal) policy.}
The gallery policy \code{refuted\_results\_must\_be\_reproved} is a pure-temporal
\emph{response} pattern (Figure~\ref{fig:pol-reproved}).
\begin{figure}[H]
\centering
\small
\[ \id{refuted\_results\_must\_be\_reproved}:\quad
\mathsf{G}\big(\id{VerificationResult}(\id{refuted}) \rightarrow \mathsf{F}(\id{VerificationResult}(\id{proved} \lor \id{sat}))\big) \]
\caption{The \code{refuted\_results\_must\_be\_reproved} policy: globally, whenever a refuted verification result appears, a proved-or-sat result must eventually follow.}
\label{fig:pol-reproved}
\end{figure}
Read: \emph{globally}, whenever a refuted verification result appears, then
\emph{eventually} a proved-or-sat result appears. Walking the semantics over the
finite timeline:
\begin{itemize}
\item \(\code{Verify}(\mathit{refuted}) \to \code{EditFile} \to
\code{Verify}(\mathit{proved})\): \textbf{passed}. At the position where the
refuted result appears, the \code{F} obligation is discharged by the later
\code{proved}.
\item \(\code{Verify}(\mathit{refuted})\) with no later proved/sat result in
scope: \textbf{failed}. The \code{G} triggers, the \code{F} obligation is never
met on the finite trace, and the evaluator reports the refuted result's action as
the witness. Note that on a finite trace \code{F} is not deferred to an infinite
future --- an obligation still open at the last position is a violation, which is
what makes ``fix by editing around it'' insufficient: only a re-established proof
closes the obligation.
\item A trace with \emph{no} refutation at all: \textbf{passed vacuously}. The
\code{G} antecedent never fires, so every position satisfies the implication.
(Depending on scope, the evaluator may instead mark this \code{NotApplicable}
when the policy's trigger genuinely never occurs; both agree that no violation
exists.)
\end{itemize}
This is the ``quality-of-evidence'' check behind the governed axis: a criterion
is \code{met} the moment a \code{VerificationResult} exists for its component, but
the result must actually \emph{establish} the property to count as \code{governed}.

\paragraph{A gallery reasoning policy.}
\code{async\_state\_transitions\_verified} governs async UI state machines
(Figure~\ref{fig:pol-async}).
\begin{figure}[H]
\centering
\small
\[ \id{async\_state\_transitions\_verified}:\quad
\mathsf{G}\big(\id{Formalize} \land \id{async\_state\_machine} \rightarrow \mathsf{F}(\id{Verify} \land (\id{proved} \lor \id{sat}))\big) \]
\caption{The \code{async\_state\_transitions\_verified} policy: a response pattern with a structural conjunct requiring every formalized async state machine to be eventually verified.}
\label{fig:pol-async}
\end{figure}
Whenever an async state machine (idle \(\to\) loading \(\to\) success/error) is
formalized, its transition logic --- especially error, timeout, and retry paths
--- must eventually be verified. It \emph{passes} when a \code{Formalize} of a
fetch machine is followed by a \code{Verify} showing all paths reach a terminal
state; it \emph{fails} for a machine where a retry-after-timeout can loop forever
(a liveness violation left unverified). This is the same response shape as above,
now with a structural conjunct in the antecedent.

\paragraph{A quantifier over the residual surface.}
Structural policies quantify over the trace's declared negative space. The
approval-gate policy \code{no\_open\_critical\_residuals} is shown in
Figure~\ref{fig:pol-critical}.
\begin{figure}[H]
\centering
\small
\[ \id{no\_open\_critical\_residuals}:\quad
\lnot\,\big(\exists\, r \in \id{residuals}.\; r.\id{severity} = \id{Critical} \land r.\id{status} = \id{Open}\big) \]
\caption{The \code{no\_open\_critical\_residuals} approval-gate policy: an existential over the residual surface asserting that no residual is both \code{Critical} and \code{Open}.}
\label{fig:pol-critical}
\end{figure}
It \emph{fails} exactly when some residual is both \code{Critical} and
\code{Open}; it \emph{passes} when the only residuals are lower-severity, or when
the critical one is \code{Addressed}/\code{Waived}. The empty-collection
semantics matter here: a trace with \emph{no} residuals passes --- \(\exists\)
over the empty surface has no witness, so its negation is true --- which is the
correct reading of ``nothing critical is open.'' The dual universal form is shown
in Figure~\ref{fig:pol-universal}.
\begin{figure}[H]
\centering
\small
\[ \id{unverified\_residuals\_have\_suggested\_check}:\quad
\forall\, r \in \id{residuals}.\; r.\id{kind} = \id{Unverified} \rightarrow r.\id{suggested\_check} \neq \emptyset \]
\caption{The dual universal form (\code{unverified\_residuals\_have\_suggested\_check}): every unverified residual must carry a suggested check, turning the negative space into a work-list.}
\label{fig:pol-universal}
\end{figure}
This form (\code{unverified\_residuals\_have\_suggested\_check}) is likewise vacuously
\emph{passed} on a residual-free trace, and turns the negative space into a
work-list on any trace that declares unverified gaps. These forms compose freely
under the temporal operators, e.g.\ the commit gate
\(\code{G}(\code{GitCommit} \to \forall r \in \code{residuals}.\,
r.\code{severity} \ge \code{High} \to r.\code{status} \neq \code{Open})\).

\subsection{Met versus governed}\label{sec:pol-met-vs-governed}

The two verdicts are \textbf{orthogonal by construction} --- how the work was
done cannot change what was established --- because they range over different
sorts of the one logic: \code{met} over the artifact sort (outcome), \code{governed}
over the action sort (process). A goal can be:
\begin{itemize}
\item \emph{met and governed} --- the evidence exists and was produced compliantly;
\item \emph{met but not governed} --- the required \code{VerificationResult}
exists (so the criterion is \code{met}), but a mandatory step was skipped or the
result was reached improperly: it was verified before being formalized, or a
refutation was edited around rather than re-proved. Rigor is a \emph{policy
decision}, not baked into recording: a verified model with no tests and no
conformance is a complete, valid artifact that resolves \code{met} by proof
alone; ``you must test/conform here'' is a pack toggle over
\code{high\_stakes\_paths};
\item \emph{governed but not met} --- every applicable process constraint holds,
but a required outcome obligation is still open.
\end{itemize}
This orthogonality is machine-checked: Table~\ref{tab:proofs} discharges
\emph{governed \(\perp\) met} alongside the goal-resolution properties (\emph{met
\(\Leftrightarrow\) all-done}, \emph{at\_risk never demotes}). And as with goals,
composition (\S\ref{sec:merge}) re-runs the policies over the \emph{merged}
timeline: a merge does not only ask ``which proofs survived,'' but ``was the
definition of done still reached in a governed way.''

\paragraph{The temporal evaluator is machine-checked.}
The soundness of this axis is not assumed. The pure decision core of the temporal
evaluator is modelled specification-first in IML and proved in ImandraX: the
\emph{policy (temporal)} model in Table~\ref{tab:proofs} discharges the
\code{G}/\code{F} semantics and \textbf{response-pattern soundness} --- the very
pattern \(\code{G}(\varphi \to \code{F}\,\psi)\) that
\code{refuted\_results\_must\_be\_reproved} instantiates --- across 78 proof
obligations, with the \emph{verdict \(\to\) residual} totality model ensuring
every terminal evaluation state lands somewhere. So when the evaluator reports a
finite-trace response policy as passed, failed, or vacuously satisfied, that
judgement conforms to a proved specification: the governance surface is itself
verified by the system it governs.


\section{Composition}\label{sec:composition}

Goals and policies judge a \emph{single} trace; \emph{composition} is what happens when traces are combined
--- at a \code{git} merge, across agents, or across reasoners. It is the sharpest test of the whole
construction, and an important technical contribution of the paper.

\subsection{Composition over traces}\label{sec:merge}

The sharpest test of ``the trace is a context'' is whether contexts \emph{compose}. The everyday instance
is a \code{git} merge; the general instance is combining two independently-produced traces (from a hub, or
from different agents or domains). A merge is the point at which two \textbf{individually-valid} results can
become \textbf{jointly invalid}: branch $A$ proves $f$ (which calls a helper $g$); branch $B$ changes $g$;
the two merge with \emph{no textual conflict and no rule broken}, yet $A$'s proof of $f$ is now stale even
though nothing in $A$'s diff changed. This cross-cutting invalidation is exactly what testing and review
miss, and it is the reason composition must be defined, not assumed.

\paragraph{Principle: merge the contract, re-derive the evidence.} The \emph{goal} (the definition of done)
is stable and mergeable; the \emph{evidence} is perishable and must be recomputed against the merged code.
Naively merging two evidence traces would \emph{launder stale-green} --- carry a proof forward as valid when
the combination invalidated it. So evidence is never merged; it is re-established, and the merge's job is to
decide, soundly, which results need re-establishing.

\paragraph{The merge operator and the affected set.}

We add \code{merge}$(\Gamma_{\text{base}}, \Gamma_{\text{ours}}, \Gamma_{\text{theirs}})$ as a first-class
trace operation --- the composition operator of the evidence logic. It records a two-parent
\code{merge\_event} (kept off the artifact DAG, so I2 acyclicity is preserved) and, for each carried-over
result $R$, computes
\[
\text{touched}(R) \;=\; \text{closure}(R) \,\cap\, \Delta,
\]
where $\Delta$ is the merge's component-wise change set and $\text{closure}(R)$ is $R$'s dependency closure
(the sound freshness signal of I4, evaluated at a two-parent point rather than against the current source).
The decision is three-way (Figure~\ref{fig:core-merge-classify}).
\begin{figure}[H]
\centering
\[
\id{classify}(R,\Delta) \;=\;
\begin{cases}
\id{CarriedForward} & \text{if } \id{touched}(R,\Delta) = \emptyset \quad\text{(closure-disjoint)}\\[2pt]
\id{CarriedForward} & \text{if } \forall d \in \id{touched}(R,\Delta).\ \id{uninterpreted}(d) \quad\text{(contract)}\\[2pt]
\id{NeedsRereasoning} & \text{otherwise}
\end{cases}
\]
where\; $\id{touched}(R,\Delta) = \id{closure}(R) \cap \Delta$.
\caption{The three-way merge classification of a carried-over result $R$ against a change set $\Delta$: carried forward when its dependency closure is disjoint from the change set, carried forward when every touched dependency was abstracted uninterpreted (a contract), and otherwise flagged for re-reasoning.}
\label{fig:core-merge-classify}
\end{figure}
The whole operator wraps this decision in a single pass over the carried-over results
(Figure~\ref{fig:merge-pseudocode}): compute the change set once, then classify each prior result and emit
its outcome, so that when the pass finishes every prior result carries exactly one verdict.
\begin{figure}[H]
\begin{lstlisting}[language=Pseudo]
merge(base, ours, theirs):
  D      <- changeset(base, ours, theirs)  # per-component changes, both sides
  merged <- combine(ours, theirs)          # union of both parents' DAGs + logs
  record two-parent merge_event(base, ours, theirs) into merged
  for each carried-over result R in merged:
    touched <- closure(R) intersect D      # I4 freshness closure at merge point
    if touched is empty:                   # closure-disjoint
      emit CarriedForward(R, witness = closure(R))
    else if every d in touched is uninterpreted:  # opaque contract
      emit CarriedForward(R, witness = touched)
    else:
      emit NeedsRereasoning(R, affected = touched)
  for each goal g with an unproven in-scope component:  # coverage layer
    regress g's obligation
  return merged   # totality: each prior result is exactly one of
                  #   CarriedForward | NeedsRereasoning | Re-reasoned
\end{lstlisting}
\caption{The \code{merge} operator in pseudocode: the change set $\Delta$ (\code{D}) is computed once, then
each carried-over result is classified (Figure~\ref{fig:core-merge-classify}) and its outcome emitted with a
witness; the coverage layer regresses a goal whose in-scope component is left unproven. On return, every
prior result is exactly one of \code{CarriedForward}, \code{NeedsRereasoning}, or (after re-establishment)
re-reasoned --- the \textbf{totality} invariant machine-checked in \S\ref{sec:proved}.}
\label{fig:merge-pseudocode}
\end{figure}
The load-bearing property is \textbf{no-false-fresh}: a carry-forward never hides a genuinely-affected
dependency. This is the freshness closure generalized, and its soundness is the \emph{same} property proved
for freshness (a fresh aggregate has no drifted contributor), so re-reasoning is bounded to the true
intersection rather than the whole trace. There are two layers: the \emph{artifact} layer (a proof's
closure drifted) and the \emph{coverage} layer (a merge adds an in-scope but unproven component --- a goal's
obligation regresses even though no existing proof changed).

\paragraph{Carried-forward as a positive claim.} Where a residual records negative space, a
\code{CarriedForward} artifact records a claim that was \emph{deliberately not re-checked because it is
provably safe}: it carries the closure it checked and a disjointness witness, so an auditor can
\emph{re-verify the skip} by recomputing the intersection. The merged trace's account of each prior result
is thus one of three, under a \textbf{totality} invariant --- \emph{carried forward}, \emph{needs
re-reasoning}, or \emph{re-reasoned}: exactly one, checkably, for every prior result. This is what makes
``we only re-reasoned the parts the merge affected'' auditable rather than a claim to be trusted.

\paragraph{Opaque contracts.} A dependency the proof \emph{abstracted} is handled by its recorded abstraction
kind. If $R$ assumed $g$ \emph{uninterpreted} --- proved its property for \emph{all} values of $g$ --- then a
merge change to $g$ cannot invalidate $R$, and $R$ carries forward (the \emph{contract} basis above). A
\emph{pinned} dependency, or a discharged contract whose callee changed, re-reasons. The abstraction
boundary is already part of the certificate (\S\ref{sec:context}); composition just reads it.

\subsection{Durable component identity}\label{sec:identity}

Composition is precise only if the change set $\Delta$ and the closures are computed over \emph{components},
not text. Binding evidence to code by name breaks under a rename: a proof of \code{clamp} stops rooting a
goal once \code{clamp} becomes \code{clamp\_int}. We give each code component a durable
\code{component\_id}, an identity layer over the descriptors an artifact already carries (per-function
hashes, target symbols, contributing sources). It sits beside the record identity \code{artifact\_id} as
\emph{which code element} versus \emph{which record}; ``the component'' is the derived group over
\code{component\_id}, and corrections are append-only aliases.

Assignment resolves, most-confident-first: shared lineage $\to$ a unique content-fingerprint match $\to$ a
confident similarity match $\to$ else \emph{mint a fresh id}. The dangerous error is \emph{conflation}
(reusing an id for a different element), so the safe fallback is to mint --- the mirror of never-guessing a
rename. With this, a goal criterion roots \emph{across} a rename, freshness \emph{follows} a rename, and the
merge change set matches a rename as one component rather than a delete plus an add.

\subsection{Generalization: multi-reasoner composition}\label{sec:multireasoner}

The evidence logic composes \emph{claims, not logics}. Each object-level reasoner emits an artifact carrying
a verdict and an explicit assumption boundary, mapped onto a shared verdict lattice
(\code{proved}/\code{refuted}/\code{unknown}/\code{bounded}/\code{instance}); \code{ponens} conjoins results
by lineage and connectives. This makes the object-logic slot pluggable: functional correctness to ImandraX,
protocols and state machines to temporal logic and model checking, arithmetic to SMT, time-and-events to the
event calculus. No unified super-logic is required --- the intractable path --- because composition happens
in the evidence logic over certified claims.

The key observation is that this is \emph{the same construction} as \S\ref{sec:merge}. A cross-branch
dependency (``$A$'s proof used $g$, which $B$ changed'') and a cross-reasoner dependency (``the arithmetic
result used a fact the temporal reasoner established'') are the same object: an \emph{explicit tracked
assumption} that must be discharged, imported, or re-checked. The soundness discipline is identical ---
every such dependency is a recorded assumption, so no implicit assumption leakage can make a conjunction
unsound --- and the trace's certificate metadata (oracle, assumptions, freshness, contested) is exactly what
makes the handoff between reasoners admissible. Composition over merges and composition over logics are one
mechanism; that is the deepest reason the trace-as-context coheres.

\part{Implementation}

\section{Formal model and verification}\label{sec:proved}

The invariants above are not informal. We modelled the framework \emph{specification-first} in IML and
proved it with ImandraX --- the system verifying its own logic; the proofs are the conformance specification
the implementation is written against. Crucially the model is not a bag of separate lemmas: it is a single
\textbf{state-transition system whose state is the trace} --- the machine-checked form of the trace-as-state
of \S\ref{sec:state}.

\paragraph{The machine.} The model is a transition system $M = (\Sigma,\ \sigma_0,\ {\rightarrow},\ Q)$. The
\textbf{state} $\Sigma$ is the trace of \S\ref{sec:state} in both its parts --- an ordered \emph{action} log
and an \emph{artifact} lineage DAG, $\sigma = (\mathit{actions}, \mathit{artifacts})$, with every artifact
produced by a recorded action (events are first-class, not a detached side-list); the \textbf{initial state}
$\sigma_0$ is empty; the \textbf{transitions} ${\rightarrow}$ are \code{extend} (record an action and the
artifact it produces), \code{supersede} (retire a target's current revision), and the two-parent
\code{combine} (a merge); and the \textbf{queries} $Q$ read the state without changing it --- freshness, goal
resolution (the \code{met} axis), policy evaluation (the \code{governed} axis), the merge affected-set
classifier, and component-identity resolution. Well-formedness $\code{wf\_state}$ --- the artifact DAG is
$\code{wf} = \mathrm{I1} \wedge \mathrm{I2} \wedge \mathrm{I3}$ (append-only, lineage-ordered/acyclic,
evidence-grounded) \emph{and} every artifact is grounded in a recorded action --- is proved an
\textbf{inductive invariant}: it holds initially and is preserved by every transition,
\[
  \code{wf\_state}(\sigma_0)
  \qquad\text{and}\qquad
  \code{wf\_state}(\sigma)\ \wedge\ \sigma \rightarrow \sigma' \ \Longrightarrow\ \code{wf\_state}(\sigma'),
\]
so \emph{every reachable state is well-formed}. The rest of the invariants are properties of the queries
over reachable states --- freshness is \emph{no-false-fresh}, the two axes are orthogonal
($\code{governed} \perp \code{met}$), and merge's carry-forward is sound --- and of the two remaining
transitions (\code{supersede} and \code{combine} also preserve \code{wf\_state}). Everything is discharged over
that one state (Table~\ref{tab:proofs}), not as a scattering of separate results.

\begin{figure}[H]
\begin{lstlisting}[language=IML]
type artifact =                                   (* a typed evidence node ... *)
  { id : int; derived_from : int list; target : int; superseded : bool;
    akind : kind; engine_backed : bool; frozen : int;
    producer_action : int }                       (* ... produced by this action *)
type action = { aid : int; atype : act; produces : int list }
type state  = { actions : action list; artifacts : artifact list }  (* THE STATE *)

let rec lineage_ordered (t : artifact list) : bool =  (* I2: edges point at earlier ids *)
  match t with [] -> true
  | a :: r -> all_earlier a.id a.derived_from && lineage_ordered r
let rec grounded (t : artifact list) : bool =         (* I3: a VResult is engine-backed *)
  match t with [] -> true
  | a :: r -> (a.akind = VResult ==> a.engine_backed) && grounded r
let wf (t : artifact list) : bool = lineage_ordered t && grounded t
(* the full-state invariant: DAG well-formed AND every artifact has a producer action *)
let wf_state (s : state) : bool =
  wf s.artifacts && artifacts_have_producers s.artifacts s.actions

(* transitions grow actions and artifacts together *)
let extend_state (s : state) (act0 : action) (a : artifact) : state =
  { actions = s.actions @ [act0]; artifacts = s.artifacts @ [a] }
let combine_state (x : state) (y : state) : state =
  { actions = x.actions @ y.actions; artifacts = x.artifacts @ y.artifacts }
(* supersede_state retires a target's current revision (flags only; log untouched) *)
\end{lstlisting}
\caption{The state and its transitions, from \code{machine/trace\_machine.iml}: the state is
\code{(actions, artifacts)} --- events are first-class, every artifact carries its \code{producer\_action};
\code{wf\_state} is the invariant; \code{extend\_state}/\code{supersede\_state}/\code{combine\_state} are the
transitions. Freshness, goals, policies, merge-classify and component identity are queries/decisions
over this same state.}
\label{fig:machine-state}
\end{figure}

\begin{table}[h]
\centering
\small
\begin{tabular}{@{}l p{9cm}@{}}
\toprule
Concern & Property proved (all over the one \code{(actions, artifacts)} state) \\
\midrule
state (events)         & \code{wf\_state}: every artifact grounded in a recorded action; preserved by \code{extend\_state}/\code{supersede\_state}/\code{combine\_state} \\
I1 append-only         & \code{wf} preserved: \code{extend} grows, \code{supersede} flag-flips \\
I2 lineage / acyclic   & \code{lineage\_ordered} preserved; no self-reference \\
I3 grounded            & prose is inert; \code{grounded} preserved \\
I4 freshness           & \code{freshness\_of} recomputed vs the current model; \textbf{no-false-fresh} \\
I5 reuse               & never-reuse-stale; conditional growth; preserves \code{wf} \\
goals (met)            & met $\Leftrightarrow$ all-done; at\_risk never demotes; progress $\in[0,1]$ \\
policies (governed)    & LTLf $G$/$F$; \textbf{governed $\perp$ met} \\
merge (composition)    & classify totality / no-false-fresh; \textbf{\code{combine} preserves \code{wf}} \\
component identity      & \textbf{never conflate}; append-only alias equivalence \\
verify escalation      & always decides; a verdict has a witness; first-decider-wins \\
rename ambiguity       & never guess when ambiguous \\
verdict totality       & every terminal verdict lands somewhere (defect $\Rightarrow$ residual) \\
\bottomrule
\end{tabular}
\caption{The whole framework proved as \emph{one} state-transition machine (the trace as state): every
invariant discharged over a single state --- \textbf{195 proof obligations, 0 failures} (ImandraX).
Composition (\S\ref{sec:merge}) and component identity (\S\ref{sec:identity}) are proved before the
implementation, so the code conforms to a proved specification.}
\label{tab:proofs}
\end{table}

The properties are proved as ImandraX \emph{verification goals} over the definitions above; a representative
selection is shown in Figure~\ref{fig:machine-goals}.

\begin{figure}[H]
\begin{lstlisting}[language=IML]
(* safety: wf_state holds initially and is preserved by every transition *)
theorem wf_state_empty = wf_state empty_state
theorem extend_state_preserves_wf (s : state) (act0 : action) (a : artifact) =
  wf_state s && admissible_state s act0 a ==> wf_state (extend_state s act0 a)
theorem combine_state_preserves_wf (ours : state) (theirs : state) =
  wf_state ours && wf_state theirs ==> wf_state (combine_state ours theirs)

(* I4 no-false-fresh: a fresh aggregate hides no drifted / vanished contributor *)
theorem no_false_fresh (t : trace) (present : int -> bool) (cur : int -> int) (a : artifact) =
  evidence_fresh t present cur && mem_art a t && is_live_result a
  ==> present a.target && a.frozen = cur a.target

(* governed _|_ met: the two axes read disjoint projections of the state *)
theorem met_ignores_timeline (s : gstate) (g : gitem list) (tl' : timeline) =
  met_of s g = met_of { s with tl = tl' } g
theorem governed_ignores_arts (s : gstate) (trig resp : act -> bool) (arts' : trace) =
  governed_of s trig resp = governed_of { s with arts = arts' } trig resp

(* merge carry-forward is sound; a component reuse is never a guess *)
theorem merge_no_false_fresh (a : artifact) (delta : int list) (u : int -> bool) =
  classify a delta u = CarriedForward
  ==> touched a delta = [] || all_uninterp u (touched a delta)
theorem reuse_is_justified (lin : int option) (ec eid b s sid id : int) =
  resolve_component lin ec eid b s sid = ReuseId id ==>
    lin = Some id || (lin = None && ec = 1 && id = eid)
                  || (lin = None && ec <> 1 && confident_sim b s && id = sid)
\end{lstlisting}
\caption{A representative selection of the verification goals \emph{proved} over the machine (all discharged
by ImandraX, 0 failures). \code{wf\_empty}/\allowbreak\code{extend\_preserves\_wf}/\allowbreak\code{combine\_preserves\_wf} are the
safety (inductive-invariant) goals; \code{no\_false\_fresh}, \code{governed $\perp$ met}
(\code{met\_ignores\_timeline} $+$ \code{governed\_ignores\_arts}), \code{merge\_no\_false\_fresh}, and
\code{reuse\_is\_justified} (never-conflate) are the flagship soundness goals.}
\label{fig:machine-goals}
\end{figure}

Three properties are worth isolating. \emph{no-false-fresh} (the merge affected-set) is the soundness of
composition: a carry-forward provably does not hide an affected dependency. \emph{never-conflate} (component
identity) is the soundness of the rebinding that makes composition rename-robust. And \emph{governed $\perp$
met} --- the two axes read disjoint projections of the state, so how the work was done cannot change what
was established --- is proved directly on the machine. That the system can express and discharge these about
\emph{itself}, over one state, is evidence that the trace-as-context is a real logical object, not a
metaphor.

\section{The Ponens system}\label{sec:impl}

The substrate is \code{ponens}, a reasoner-agnostic trace format and toolchain: an append-only typed trace,
a residual surface, a goal/policy layer, and derived freshness, with a CLI (\code{trace enrich},
\code{trace check}, and \code{trace merge}). Merge is realized as described in \S\ref{sec:merge}: a report
mode and a \code{--combine} mode that emits a materialized merged trace which the validator, enricher, and
policy checker accept. The affected-set classifier reuses the closure-checksum machinery that already
computes freshness, so composition is freshness ``evaluated at a two-parent point.'' The format spec is
versioned and additive; the composition additions (the two residual kinds, the carried-forward artifact, the
\code{component\_id} field, and the merge section) were added as a backward-compatible version bump.

The concrete instantiation is a coding agent whose object logic is IML/ImandraX. Here composition is
triggered on real developer workflow: a \code{git} post-merge hook recognizes a merge (distinct from an
ordinary edit), and after the freshness recompute surfaces exactly which verified results the merge put at
risk --- ``this merge put $N$ verified results at risk; re-verify to restore the guarantee.'' The default is
to re-enrich and flag; re-verification is an explicit, bounded follow-up over the at-risk set. The whole
addition is additive: a trace without composition metadata behaves exactly as before.

\subsection{The CLI in practice}\label{sec:impl-cli}

Every construct in this paper is exercised through one command-line tool, \code{ponens}, operating on a local
trace file; the same verbs run inside an agent's harness and in a reviewer's terminal. A representative
lifecycle begins by deriving a trace from an agent session, then declaring the definition of done, earning
its resolution against the recorded evidence, and governing \emph{how} the work was produced
(Figure~\ref{fig:core-cli-lifecycle}).

\begin{figure}[H]
\begin{lstlisting}[language=PonensCLI]
# Derive a trace from the latest agent session transcript (or `ponens trace init`)
ponens emit -o trace.json --summarize

# Enrich: resolve acceptance, derive the residual surface and relevance cones,
# and recompute per-result freshness (Fresh | Stale | Detached)
ponens trace enrich trace.json

# Goals: state intent + definition of done, bind an acceptance item to a symbol
ponens trace goal set    trace.json --intent "no double-charge on retry" --scope pricing.py
ponens trace goal accept trace.json --kind property --label "idempotent charge" --symbol charge

# Resolve the goal against the evidence (earned, not self-reported) ...
ponens trace resolve trace.json
# ... and record a non-doer's sign-off that the definition of done is RIGHT
ponens trace goal certify trace.json --by reviewer --verdict approved

# Govern: load best-practice policies and run the gate (fully offline)
ponens registry update
ponens policies add tests_before_commit --into trace.json
ponens trace check trace.json               # Passed / Failed / NotApplicable, per policy
ponens trace report trace.json              # Markdown grade + open gaps for a PR comment
\end{lstlisting}
\caption{A representative \code{ponens} lifecycle on a local trace file: deriving a trace from an agent session, enriching it, declaring the goal and definition of done, resolving it against the recorded evidence, and governing how the work was produced with best-practice policies.}
\label{fig:core-cli-lifecycle}
\end{figure}

The interchange and non-repudiation verbs (\S\ref{sec:interchange}) sign the frozen content, bind it 1:1 to
the commit under review, and project the lineage to a standard provenance vocabulary (Figure~\ref{fig:core-cli-interchange}).

\begin{figure}[H]
\begin{lstlisting}[language=PonensCLI]
# Sign the content_hash with a role-bearing disposition; verify against a roster
ponens trace sign   trace.json --role auditor --disposition approved   # ssh | gpg | sigstore
ponens trace verify trace.json --allowed-signers roster

# Bind to the git commit (a Trace-Id git note) and publish for review
ponens bind && ponens push

# Project the typed-artifact lineage to W3C PROV-JSON for an external auditor
ponens trace export trace.json --to prov -o trace.prov.json
\end{lstlisting}
\caption{The interchange and non-repudiation verbs: signing the frozen \code{content\_hash} with a role-bearing disposition and verifying it against a roster, binding the trace to the git commit under review, and exporting the typed-artifact lineage to W3C PROV-JSON for an external auditor.}
\label{fig:core-cli-interchange}
\end{figure}

Composition (\S\ref{sec:merge}) is a single verb. In report mode it names the verified results a merge put
at risk; \code{--combine} materializes the merged trace, which the enricher, checker, and validator accept
unchanged (Figure~\ref{fig:core-cli-merge}).

\begin{figure}[H]
\begin{lstlisting}[language=PonensCLI]
# Merge two branch traces: carry forward the provably-unaffected, flag the rest
ponens trace merge ours.json theirs.json --base base.json --combine -o merged.json
ponens trace enrich merged.json             # re-derive freshness at the two-parent point
\end{lstlisting}
\caption{Composition as a single verb: \code{ponens trace merge} carries forward the provably-unaffected results and flags the rest, and \code{--combine} materializes a merged trace whose freshness is then re-derived at the two-parent point.}
\label{fig:core-cli-merge}
\end{figure}

\subsection{The LTL policy checker}\label{sec:impl-ltl}

The governed axis of \S\ref{sec:policies} is enforced by a concrete checker built into \code{ponens}. A
policy formula --- written in the unicode surface of \S\ref{sec:policies} or an equivalent ASCII surface ---
is tokenized and parsed by a recursive-descent parser into an abstract syntax tree over the operators of the
language: the future-time \code{G}/\code{F}/\code{X} and the binary \code{U}; the past-time
\code{P}/\code{H}/\code{S}/\code{S\_last}; the scoped-past \code{P\_target}/\code{P\_chain}/\code{H\_target}/%
\code{H\_chain}; the first-order quantifiers $\forall$/$\exists$/$\exists!$ over the typed artifact graph;
field access and comparison; set comprehensions; and the \code{count} aggregate. The evaluator interprets
that AST over the finite trace --- the ordered action timeline --- and returns a three-valued verdict per
policy: \code{Passed}, \code{Failed}, or \code{NotApplicable} (the trigger never fired), together with
\emph{witnesses}: the action and artifact ids that support a pass or break a fail
(Figure~\ref{fig:ltl-checker}). \code{ponens trace check} runs every attached policy this way and aggregates
the verdicts into the audit grade (or, for the FIX pack of \S\ref{sec:examples-fix}, the Green/Amber/Red
\code{GovernanceState}).

\begin{figure}[H]
\begin{lstlisting}[language=Policy]
# a policy formula ...            parse -> AST -> evaluate over the action timeline
G( Release -> P(UserApproval /\ authenticated) )
\end{lstlisting}
\vspace{-6pt}
\begin{lstlisting}[language=PonensCLI]
$ ponens trace check trace.json
  no_release_without_authenticated_approval  FAIL  (Release@a17: no prior auth UserApproval)
  tests_before_commit                        PASS
  review_with_independence                    N/A  (trigger never fired)
  grade: AMBER
\end{lstlisting}
\caption{The LTL checker (schematic): a formula is parsed to an AST and evaluated over the action timeline,
yielding \code{Passed}/\code{Failed}/\code{NotApplicable} per policy with witnessing ids, which
\code{ponens trace check} aggregates into a grade.}
\label{fig:ltl-checker}
\end{figure}

Two things anchor the checker's trustworthiness. First, the check is \textbf{deterministic and fully
offline} --- no model in the loop --- which is exactly the property a \emph{binding} governance control must
have (\S\ref{sec:policies}; the FIX proposal's ``deterministic function over the governed state'' of
\S\ref{sec:examples-fix}). Second, the language has \emph{two evaluators kept in lockstep} --- the Python CLI
evaluator and a JavaScript evaluator (the browser playground and trace viewer) --- checked against each
other by a cross-evaluator parity harness, so a policy yields the same verdict wherever it runs. And the very
same formulas \emph{compile to IML} (over a policy-combinator library), so a policy can be handed to
ImandraX; the governed axis itself is machine-checked in the state machine of \S\ref{sec:proved} (the
\textrm{LTL}$_f$ \code{G}/\code{F} semantics and \code{governed} $\perp$ \code{met}). The runtime checker and
the formal model are two readings of one language.

\section{Interchange and sync}\label{sec:interchange}

Two more outward faces remain, both about \emph{transport} rather than trust: how the record is read by
tooling that was not built for it, and how it travels with the code it explains. Both are deliberately
narrow. A provenance projection carries lineage and attribution and nothing else; a git binding carries the
trace's identity and nothing else. What the record \emph{means} --- freshness, contested-ness, the residual
surface, the verdicts --- stays with \code{ponens}, where it can be recomputed.

\subsection{Interchange: projecting to W3C PROV}\label{sec:interchange-prov}

An agent's provenance ought to be legible to tools that were not built for our format. We therefore project
\trace{} onto \textbf{W3C PROV}~\cite{prov}, the standard vocabulary for provenance
(\code{ponens trace export --to prov}, emitting PROV-JSON). The mapping is close to structural: a typed
artifact becomes a \code{prov:Entity}, an action a \code{prov:Activity}, the assistant (with its model) a
\code{prov:Agent}, and the trace itself a \code{prov:Bundle}. The lineage edges carry over one-to-one ---
\code{producer\_action\_id} to \code{wasGeneratedBy}, an action's inputs to \code{used},
\code{derived\_from} to \code{wasDerivedFrom}, and \code{supersedes} to a \code{wasDerivedFrom} of type
\code{prov:Revision} (i.e.\ \code{wasRevisionOf}). Every artifact and action is \code{wasAttributedTo} /
\code{wasAssociatedWith} the agent. Edges whose endpoints are not present in the trace (a dangling
\code{derived\_from}) are dropped rather than emitted, so the result is a well-formed PROV instance.

The essential point --- and the reason this belongs beside \S\ref{sec:context} rather than in place of it ---
is that \emph{PROV is an interchange view, not a reasoning state}. PROV models \emph{what was derived from
what, by whom}; it has no vocabulary for the evidence-logic certificate that makes the trace a context in the
first place. Concretely, the projection is \textbf{faithful for lineage and attribution} and \textbf{lossy for
judgement}:

\begin{itemize}
\item \textbf{Freshness} (I4) is \emph{derived} by \code{ponens} from a dependency-closure fingerprint; PROV's
\code{wasInvalidatedBy} records an invalidation \emph{event someone asserts}, not a \emph{computed} staleness,
so it is not emitted.
\item \textbf{Residuals and defeaters} --- the negative space and the contested-ness of a claim --- have no
PROV relation. A \code{Residual} is exported as a \code{prov:Entity} with its \code{kind}, \code{severity},
and \code{status} riding as \code{ponens:} attributes, but PROV assigns them \emph{no meaning}; the blocking of
a contested claim by counter-evidence stays a \code{ponens}-side computation.
\item \textbf{Verification verdicts} (proved / refuted / unknown) ride as an opaque \code{ponens:verdict}
attribute on the entity.
\item \textbf{Policies} --- the temporal-logic governance axis --- have no analogue; PROV-CONSTRAINTS is a
fixed consistency checker, not a policy language.
\end{itemize}

The rule of thumb is \emph{lineage and attribution round-trip; judgement does not}. Export is deliberately
one-way: a consumer that needs the certificate --- freshness, contested-ness, the residual surface, the
policy verdict --- reads the \code{ponens} trace, and the standard vocabulary is offered so that provenance
tooling can at least consume chain-of-custody and derivation without learning our format. This is the same
distinction the paper draws throughout: a log records \emph{what happened}; a context carries \emph{the
warrant}. PROV is the log projection of a state that is more than a log.

\subsection{Sync and binding}\label{sec:interchange-sync}

For the warrant to be \emph{useful} to a downstream party, it must travel with the code it explains. Three
facts live in three places and no single layer owns all three: the \textbf{local trace file} is canonical for
\emph{content} (immutable, content-hashed); the \textbf{git commit} is canonical for \emph{version} (which
code state the reasoning is about); and a \textbf{hub} is canonical for \emph{decisions} (status, comments,
review items, snapshots, policy runs). The CLI's job is to bind them and keep them from silently diverging ---
in one line, it makes the hub a git remote for reasoning.

\code{ponens bind} is the load-bearing verb. It stamps the trace with \code{repo} / \code{branch} /
\code{commit\_sha} from git, and writes the trace identity \emph{back} onto the commit --- as a
\code{Trace-Id} trailer or, by default, a \code{git} note (\code{refs/notes/ponens}, non-mutating). This makes
the relationship resolvable in both directions, $\code{commit\_sha} \Leftrightarrow \code{trace\_id}$, so the
reasoning is discoverable from the commit (e.g.\ in a pull-request view) and vice versa. The binding is
\textbf{1:1}: a trace is bound to exactly one commit, and iteration produces a \emph{new} commit and a
\emph{successor} trace rather than overwriting the old one --- the same immutability discipline as
\S\ref{sec:interchange-review}, now mirrored one-to-one with git history. \code{push} / \code{pull} move the
trace and the hub's decision state (\code{push} is idempotent and creates a successor on a content change;
\code{pull} is read-only with respect to the working tree), and \code{status} names the one steady state
(\code{synced}) against the divergences (\code{stale}, \code{dirty}, \code{behind}) that each map to a concrete
next command.

The design decision is \textbf{binding-only}: git carries the trace and the link, and \emph{only} those. This
is what lets the warrant travel with the code \emph{without} committing the churny raw trace state or
reinventing review inside git --- git notes do not merge cleanly and committed review files conflict on every
parallel comment (this is the local regime the merge composition of \S\ref{sec:merge} operates within: git
merges the code and the binding; \code{ponens} re-derives the evidence). Human review of the \emph{diff} stays
on the git platform; computable governance and the frozen, signed \emph{snapshot} (the portable audit artifact
with no platform equivalent) stay on the hub; the commit is the shared anchor connecting them. Let the
platform own diff review, let the hub own governance over the reasoning, and connect them through the commit
--- do not reimplement either inside the other.



\part{Demonstration and discussion}

\section{Worked examples}\label{sec:examples}

The construction is easiest to see end-to-end. We walk two scenarios --- one from software development, one
from agentic trading --- each exercising the full machinery: a goal, its resolution against evidence,
governance by policy, and (for the first) composition across a merge. The two differ only in their
\emph{object logic} and \emph{policy vocabulary}; the evidence logic over the trace is identical. The trading
example runs the FIX Trading Community's runtime-governance pack~\cite{ponensfix} verbatim.

\subsection{A verified pricing change through a merge}\label{sec:examples-coding}

An agent is asked to make a payment function \code{charge} idempotent under retries: a retried request must
never double-charge. This pays off motivating example~1 (\S\ref{sec:intro-examples}) concretely.

\begin{enumerate}
\item \textbf{Declare the goal.} The definition of done is stated and bound to the component it constrains
(\S\ref{sec:goals}):
\begin{quote}
\code{ponens trace goal set trace.json --intent "charge is idempotent under retry" --scope pricing.py}\\
\code{ponens trace goal accept trace.json --kind property --label "no double-charge on retry" --symbol charge}
\end{quote}
The goal is now a first-class field of \trace{}; its acceptance item names the property and the component
\code{charge}, not a line of prose.

\item \textbf{Establish the evidence.} The agent formalizes \code{charge} in IML, defines a verification goal,
and ImandraX discharges it. A \code{VerificationResultArtifact} \code{vr1} with status \code{proved} enters the
lineage, carrying a \code{reasoning\_fingerprint} over \code{charge} \emph{and its dependency closure} --- which
includes the helper \code{apply\_discount} (\S\ref{sec:format}).

\item \textbf{Resolve and govern.} \code{ponens trace resolve} walks the lineage: an artifact of the required
type roots in \code{charge}, is \code{Fresh}, and is uncontested, so the acceptance item resolves \code{met}
and the goal is met --- earned from evidence, not self-reported (\S\ref{sec:goals}). \code{ponens trace check}
runs the governance pack (e.g.\ tests-before-commit) and returns Green (\S\ref{sec:policies}); \code{ponens
trace sign} and \code{ponens bind} freeze and attribute the record (\S\ref{sec:interchange-sign}).

\item \textbf{A change arrives on another branch.} A colleague, on a separate branch, refactors
\code{apply\_discount} --- a function in \code{vr1}'s dependency closure --- to change rounding. Their branch's
own trace is internally consistent and never mentions \code{charge}; in isolation, each side is valid.

\item \textbf{The merge, decided soundly.} A \code{git} merge fires the post-merge hook, which runs
\code{ponens trace merge ours.json theirs.json --base base.json --combine} (Figure~\ref{fig:core-cli-merge}).
The classifier intersects \code{vr1}'s dependency closure with the merge's change set: \code{apply\_discount}
is touched and was not abstracted \code{uninterpreted}, so \code{vr1} is \emph{not} carried forward --- it
becomes a \code{NeedsRereasoning} residual and the goal's verdict drops to \code{at\_risk} (\S\ref{sec:merge}).
The report reads: ``this merge put 1 verified result at risk; re-verify to restore the guarantee.'' By the
\emph{no-false-fresh} property (proved, \S\ref{sec:proved}), the proof is never silently kept.

\item \textbf{Close the loop.} Re-running the proof over the at-risk set either restores \code{proved} (the goal
returns to \code{met}, freshness \code{Fresh}) or produces a counterexample --- a \code{Defeater} that
\code{Rebuts} \code{vr1}, making the claim contested and resolving the acceptance item \code{AcceptBlocked}
(\S\ref{sec:goals}). Either way the result is a checkable verdict, not a silent regression: the
individually-valid proof and the individually-valid refactor were \emph{jointly} invalid, and the trace made
that decidable.
\end{enumerate}

\subsection{Governing an agentic trading workflow (FIX)}\label{sec:examples-fix}

An execution agent works an order in secondary markets. Its run is recorded as a trace whose actions use the
agentic vocabulary the FIX AI Working Group's runtime-governance proposal~\cite{fixgov} assumes ---
\code{Retrieve} (market data), \code{Compute} (order size), \code{Draft} (the order), \code{Release} (send to
venue) --- each carrying the governance predicates the runtime emits per action (identity resolved, vLEI
present, intent resolved, tool allowlisted, provenance and recency checked, deterministic recompute,
authenticated approval). We govern this trace with the FIX Trading Community's \emph{Agentic AI Runtime
Governance} pack~\cite{ponensfix}, which maps the proposal's traffic-light (Green / Amber / Red) scheme onto
\code{ponens} policies (Figure~\ref{fig:ex-fix-pack}); the underlying capability model and four-tier risk
scheme are from Szpruch et al.~\cite{szpruch}. The same shape appears in their independently-developed
credit-memo example --- a labelled transition system with a guard on every transition (retrieval recency,
deterministic recomputation, an authenticated-approval gate before release). Read in our terms, that
trajectory is a trace and those guards are policies of the form in Figure~\ref{fig:ex-fix-pack}, which
\code{ponens trace check} evaluates as the deterministic governance function both lines of work call for.

\begin{figure}[H]
\centering
\small
\[
\begin{array}{@{}l@{\;:\;}l@{}}
\id{no\_release\_without\_authenticated\_approval} & \text{\footnotesize\itshape \textcolor{red!70!black}{Red}; approval \& release gating}\\
\multicolumn{2}{@{}l@{}}{\quad \mathsf{G}\big(\id{Release} \rightarrow \mathsf{P}(\id{UserApproval} \land \id{authenticated} \land \id{approval\_scope\_covers})\big)}\\[4pt]
\id{dual\_approval\_critical} & \text{\footnotesize\itshape \textcolor{red!70!black}{Red}; Tier 4, two approvers}\\
\multicolumn{2}{@{}l@{}}{\quad \mathsf{G}\big(\id{Release} \rightarrow \mathsf{P}(\id{UserApproval} \land \id{approver\_1}) \land \mathsf{P}(\id{UserApproval} \land \id{approver\_2})\big)}\\[4pt]
\id{retrieved\_data\_attributed} & \text{\footnotesize\itshape \textcolor{red!70!black}{Red}; Szpruch C1 (retrieval)}\\
\multicolumn{2}{@{}l@{}}{\quad \mathsf{G}\big(\id{Retrieve} \rightarrow \id{provenance\_checked} \land \id{recency\_checked}\big)}\\[4pt]
\id{numeric\_recomputed\_deterministically} & \text{\footnotesize\itshape \textcolor{red!70!black}{Red}; Szpruch C2 (numerics)}\\
\multicolumn{2}{@{}l@{}}{\quad \mathsf{G}\big(\id{Compute} \rightarrow \id{deterministic\_recompute}\big)}\\[4pt]
\id{credential\_not\_expiring} & \text{\footnotesize\itshape \textcolor{orange!85!black}{Amber}; identity \& auth}\\
\multicolumn{2}{@{}l@{}}{\quad \mathsf{G}\big(\id{action} \rightarrow \lnot\,\id{credential\_expiring}\big)}
\end{array}
\]
\caption{A slice of the FIX Trading Community runtime-governance pack~\cite{ponensfix,fixgov}: traffic-light
conditions expressed as temporal-logic policies over the trading agent's execution trace. \code{ponens trace
check} aggregates the per-policy verdicts into a single Green / Amber / Red \code{GovernanceState}.}
\label{fig:ex-fix-pack}
\end{figure}

\begin{enumerate}
\item \textbf{Attach and check.} \code{ponens trace check trace.json} evaluates each traffic-light condition as
a policy over the trace and aggregates the verdicts: all-pass is Green, a warning-severity failure is Amber, an
error-severity failure is Red (\S\ref{sec:policies}).

\item \textbf{A red flag.} Suppose the agent issued a \code{Release} with no prior authenticated,
scope-covering approval. \code{no\_release\_without\_authenticated\_approval} (Figure~\ref{fig:ex-fix-pack})
fails at \code{error} severity, so the aggregate is \textbf{RED}: \code{ponens trace check} exits non-zero and
names the failed policy --- the \code{GovernanceFlags} the proposal wants carried in a new FIX field. In
parallel a credential nearing expiry trips \code{credential\_not\_expiring} at \code{warning} severity
(\textbf{Amber}), which is reported but never collapsed into the Red.

\item \textbf{Deterministic, model-independent enforcement.} The verdict is a deterministic function over the
trace, executing independently of the language model --- exactly the property the proposal requires of any
\emph{binding} (as opposed to advisory) control~\cite[\S10.3]{fixgov}: ``a governance decision that cannot be
expressed as a deterministic function over the governed state \ldots\ is advisory guidance, not binding
enforcement.'' A \code{ponens} policy \emph{is} that function.

\item \textbf{Remediate and re-gate.} A Tier-4 (Critical Autonomous) deployment additionally requires
\code{dual\_approval\_critical}. Once the agent obtains authenticated dual approval before \code{Release},
\code{ponens trace check} returns Green; the signed, bound trace (\S\ref{sec:interchange}) is the auditable
\code{GovernanceState} record --- verifiable by the venue or a regulator without trusting the agent's operator.
\end{enumerate}

In both examples the same object --- a proof-carrying trace queried by the evidence logic --- carries the
verdict; only the object logic (IML/ImandraX proofs versus governance predicates over a trading trajectory)
and the policy vocabulary differ. That is the multi-reasoner thesis of \S\ref{sec:multireasoner} in miniature.


\section{Related work}

\emph{Provenance.} W3C PROV~\cite{prov} standardizes \emph{how} an artifact came to be, for interchange; it
is a lossy view of, not a substitute for, a reasoning state --- it carries neither the evidence-logic
certificate (freshness, contested-ness) nor a notion of composition. \emph{Non-monotonic and defeasible
reasoning.} The trace is a defeasible context in Pollock's sense~\cite{pollock}: defeaters record
counter-evidence and retract conclusions; the query-carrying property connects it to answer-set
programming~\cite{asp} and the event calculus~\cite{ec}. \emph{Proof-carrying code}~\cite{pcc}: claims travel
with certificates checkable by the consumer; we extend the idea from a single artifact to a composable
context. \emph{Truth-maintenance systems}~\cite{atms}: dependency-directed retraction is the local analogue
of our non-monotonic, well-founded context --- the trace can be seen as a proof-carrying, append-only TMS
over a codebase, with merges as the composition operation TMSs do not define. \emph{Version control.} Git
infers renames heuristically at diff time and merges text without any notion of semantic invalidation;
\S\ref{sec:identity} gives the durable identity git omits, and \S\ref{sec:merge} the semantic merge it does
not attempt.

\emph{Runtime governance for agentic AI.} A parallel line in financial Model Risk Management, developed
independently of this work, arrives at strikingly similar premises: agentic risk is a property of
\emph{execution trajectories} rather than of an input--output mapping, and binding governance must be a
\emph{deterministic function over the governed state}, executed independently of the language
model~\cite{szpruch,fixgov}. That is the object and the discipline we arrive at here as well --- their
governed trajectory is our trace, and their governance function is a \code{ponens} policy
(\S\ref{sec:policies}) that \code{ponens trace check} evaluates --- and we read the convergence as mutual
corroboration rather than either deriving from the other. Ponens itself was driven by industrial applications
in capital markets --- certifying counterparty connectivity against venue FIX specifications, and governing
agentic trading workflows --- rather than by a governance theory; the convergence is therefore between field
practice and an independently-derived MRM programme, which we take as the stronger evidence that treating the
execution trajectory as the governed object is the right move. The agreement runs deeper on two points that
our construction independently makes concrete. Both hold that verification must be \emph{decomposed by claim
structure} --- logical, numeric, retrieval, consistency --- rather than delegated to one semantic judge,
which here is multi-reasoner composition (\S\ref{sec:multireasoner}); and both hold that an LLM-as-judge
conflates \emph{plausibility} with \emph{correctness}, which here is the \textsc{established}-versus-%
\textsc{asserted} discipline (\S\ref{sec:interchange-audit}). The two frameworks are complementary in scope:
theirs is a capability-catalogue MRM \emph{programme} for a bank; ours contributes a concrete trace format, a
policy language and its formal semantics, the goal axis, \emph{composition over traces at a merge} --- a
notion their setting, concerned with single-run orchestration and change control, does not address --- and
machine-checked foundations that make the governance function checkable and its warrant portable across
parties.

\section{Conclusion}

Treating the reasoning trace as a state --- a defeasible, proof-carrying context rather than a log --- is not
bookkeeping; it is what makes agent reasoning \emph{compose}. The same append-only structure serves as
memory, governance surface, and logical context, and the certificate every claim carries is exactly what a
merge (or a second reasoner) needs to decide, soundly, what a change preserves and what it invalidates. We
gave the model, its composition operator, and a durable identity for binding evidence to code across change;
we proved the load-bearing invariants by the system on itself; and we showed the construction generalizes
from merges to multiple logics. The result is a substrate on which an agent's work is not merely recorded but
\emph{verifiable} --- and, because the warrant travels with the claim, verifiable by a party other than the
one that produced it.

\paragraph{Availability.} \code{ponens} is open source. The toolchain and the full trace, goal, and policy
specifications are at \url{https://www.ponens.dev}; the source is at
\url{https://github.com/imandra-ai/ponens}.

\begin{thebibliography}{9}
\bibitem{imandra} Passmore, G.\ et al.\ \emph{The Imandra Automated Reasoning System.} IJCAR, 2020.
\bibitem{prov} Moreau, L., Missier, P.\ (eds.).\ \emph{PROV-DM: The PROV Data Model.} W3C Recommendation, 2013.
\bibitem{pollock} Pollock, J.\ \emph{Defeasible Reasoning.} Cognitive Science 11(4), 1987.
\bibitem{asp} Gelfond, M., Lifschitz, V.\ \emph{The Stable Model Semantics for Logic Programming.} ICLP, 1988.
\bibitem{ec} Kowalski, R., Sergot, M.\ \emph{A Logic-Based Calculus of Events.} New Generation Computing 4, 1986.
\bibitem{ltl} Pnueli, A.\ \emph{The Temporal Logic of Programs.} FOCS, 1977.
\bibitem{pcc} Necula, G.\ \emph{Proof-Carrying Code.} POPL, 1997.
\bibitem{atms} de Kleer, J.\ \emph{An Assumption-Based TMS.} Artificial Intelligence 28(2), 1986.
\bibitem{fixgov} FIX Trading Community, AI Working Group.\ \emph{Proposal on Agentic AI Runtime Governance under FIX Protocol Extensions.} R.\ Healey \& K.\ Houston, 12 Jun 2026.
\bibitem{szpruch} Szpruch, L., Sudjianto, A., Bhatti, T., Ang, G.\ \emph{Scalable Runtime Governance for Agentic AI in Financial Services.} SSRN 6567199, 2026.
\bibitem{ponensfix} FIX Community.\ \emph{Agentic AI Runtime Governance --- \code{ponens} policy pack.} \url{https://www.ponens.dev/organizations/fix-community}.
\end{thebibliography}

\end{document}
